% Standalone source. All manuscript text, references, and printed listings are embedded.
\documentclass[reqno]{amsart}
\usepackage[T1]{fontenc}
\usepackage{lmodern}
\usepackage{amsmath,amssymb,mathtools,mathrsfs}
\usepackage{booktabs,longtable,array}
\usepackage{enumitem}
\usepackage{xcolor}
\usepackage{microtype}
\usepackage[unicode,pdfusetitle,hidelinks]{hyperref}
\usepackage{listings}
\setlength{\emergencystretch}{2em}
\setcounter{tocdepth}{1}
\allowdisplaybreaks[2]
\newtheorem{theorem}{Theorem}[section]
\newtheorem{proposition}[theorem]{Proposition}
\newtheorem{lemma}[theorem]{Lemma}
\newtheorem{corollary}[theorem]{Corollary}
\theoremstyle{definition}
\newtheorem{definition}[theorem]{Definition}
\theoremstyle{remark}
\newtheorem{remark}[theorem]{Remark}
\newtheorem{example}[theorem]{Example}
\newcommand{\C}{\mathbb C}
\newcommand{\R}{\mathbb R}
\newcommand{\N}{\mathbb N}
\newcommand{\D}{\mathbb D}
\newcommand{\HH}{\mathbb H}
\DeclareMathOperator{\Tr}{Tr}
\DeclareMathOperator{\Id}{Id}
\DeclareMathOperator{\supp}{supp}
\DeclareMathOperator{\rank}{rank}
\DeclareMathOperator{\diag}{diag}
\DeclareMathOperator{\osc}{osc}
\DeclareMathOperator{\dist}{dist}


\title[Positive stability semigroups and capacities]{Positive stability-preserving semigroups\\ and particle capacities}
\author{Dongsheng Wei}
\address{Independent Researcher}
\email{dongshengwei2025@icloud.com}
\date{}
\subjclass[2020]{Primary 47D06, 60J27; Secondary 30C15, 60K35}
\keywords{stable polynomial, positive semigroup, Laguerre--P\'olya class, particle system, generator classification, finite capacity, semigroup generation}
\hypersetup{pdftitle={Positive stability-preserving semigroups and particle capacities},pdfauthor={Dongsheng Wei}}
\begin{document}
\begin{abstract}
We classify the positive strongly continuous semigroups on summable
coefficient spaces that preserve stable inputs with nonnegative
coefficients, allowing finite, infinite, and mixed coordinate capacities.
On the unrestricted state
space, preservation forces a canonical differential operator of order
at most two and a logarithmic symbol with finitely many parameters.
For conservative nonexplosive Markov chains, no prior restriction on
jump sizes is needed: preservation allows only constant immigration,
independent migration, and deaths with a common quadratic collective
rate. At finite and mixed capacities, an explicit rate classification
and a universal quadratic-contact inequality give an exact criterion.
Complete preservation under inert extensions is detected by two binary
ancillary coordinates, or one coordinate of capacity two, and reduces
to entrywise inequalities on the principal coefficient matrix. The
infinite-dimensional statements distinguish classification of an
existing semigroup from generation of a formal operator. An automatic
domain theorem, a consistent-cutoff criterion, and compatible weighted
spaces supply the required first variations and finite-capacity limits.
\end{abstract}
\maketitle
\tableofcontents
\section{Introduction}

Real-stable generating polynomials encode strong restrictions on a
probability distribution. The associated preservation problem asks
which dynamics respect these restrictions at every time. For binary
particle systems, stability is the defining algebraic property of
strongly Rayleigh measures. For populations with no fixed capacity,
the natural closed class consists of entire functions in the positive
Laguerre--P\'olya class. These settings share polynomial tests but differ
in their boundary constraints and in the analytic meaning of a
generator.

The theory of stable-polynomial preservers developed by Borcea and
Br\"and\'en~\cite{BB09,BB10,BB10II} gives symbol criteria for linear
operators preserving the whole stable class. Borcea, Br\"and\'en, and
Liggett~\cite{BBL09} connect stability with negative dependence and
stability-preserving particle dynamics. Liggett and
Vandenberg-Rodes~\cite{LVR10} treat stability on unbounded population
spaces, including birth--death chains and particle systems. We use
their positive Laguerre--P\'olya closure and approximation results, as
well as the established symbol and polarization machinery. Our subject
is the classification of continuous positive preservation across
coordinate capacities, with the operator domain and existence questions
included explicitly.

The positive restriction changes the infinitesimal problem. An operator
need only preserve stable polynomials with nonnegative coefficients,
and such preservation need not survive the addition of inert variables.
Collective particle loss is an example. The second-order sign conditions
for preservation of arbitrary complex stable polynomials therefore do
not classify this positive problem. Two boundary tests replace them:
collisions of finite roots eliminate canonical derivatives of order
greater than two, and first variations at Poisson generating functions
restrict the degree and signs of the surviving coefficients.

\subsection{The classification and its probabilistic specialization}

Theorem~\ref{pos:symbol} classifies existing positive
\(C_0\)-semigroups on \(\ell^1(\N^d)\) preserving the positive
Laguerre--P\'olya class. Its canonical expression consists of affine
first-order terms, multiplication by a nonnegative linear polynomial,
collective second-order loss, common quadratic damping, and individual
quadratic damping. Equivalently, its logarithmic symbol has the precise
finite-parameter form in \eqref{pos:log-symbol}. The parameters specify
an algebraic cone; bounded semigroup realization on the chosen normed
space is a separate question.

For conservative nonexplosive Markov chains, the mass identities
specialize this cone sharply. Corollary~\ref{pos:markov} starts with
arbitrary jumps and arbitrary finite state-dependent exit rates. It
shows that preservation is equivalent to constant immigration,
independent migration, and deaths
\[
 x\longrightarrow x-e_i
 \quad\text{at rate}\quad x_i[\delta_i+\gamma(|x|-1)],
 \qquad \delta_i,\gamma\ge0.
\]
The quadratic coefficient \(\gamma\) is common to all sites. It is
enough to test multinomial generating polynomials
\((c+a\cdot z)^m\), with \(c\ge0\), \(a>0\), and arbitrary \(m\).
An exact positive Euler step for collective loss provides the elementary
preserving flow used in the sufficiency proof.

\subsection{Capacities and quadratic contacts}

For a finite capacity vector \(\kappa\), Theorem~\ref{pos:finite}
first reduces the allowed transitions to explicit death, birth, and
transfer rates. Their capacity factors record the faces of the state
box. Within this rate family, preservation is equivalent to
\[
 \operatorname{tr}(Q(y)H)\le0
 \qquad ((y,H)\in\mathfrak Q_J),
\]
where \(Q\) is the symmetric principal coefficient matrix and
\(J=\{i:\kappa_i=1\}\). The contact set \(\mathfrak Q_J\), defined
in \eqref{pos:contact-domain}, has a fixed quadratic form: beyond the
binary coordinates, numerical capacities enter the rates rather than
the contact geometry. Every failure of this inequality has a quadratic
stable first-order witness. Theorem~\ref{pos:mixed} gives the corresponding
criterion when some coordinates have infinite capacity. Its proof uses
finite approximations with unchanged diagonals, so that artificial
boundaries introduce killing rather than an unverified reflecting flow.

Theorem~\ref{gen:ancilla} identifies exactly which positive semigroups
remain preserving after inert extensions. Two new binary variables, or
one variable of capacity two, suffice to test every finite collection
of ancillary capacities, including infinite ones. The criterion is
\(Q_{ij}(y)\le0\) for every real \(y\) and every pair of indices.
One inert binary variable does not suffice uniformly. In a finite box,
the resulting cone agrees with complete preservation of arbitrary
complex stable polynomials; this final comparison uses the real
generator theorem from the companion manuscript~\cite{WeiGenerators2026}.
The positive classifications and the ancillary test themselves are
proved here.

\subsection{Domains, generation, and organization}

On an infinite state space, writing a formal differential expression
does not specify a bounded-operator \(C_0\)-semigroup. Conversely, a
given positive \(C_0\)-semigroup need not initially come with
polynomials in its generator domain. Theorem~\ref{ana:domain} supplies
this domain conclusion from log-concavity of the graded transition
columns, using the countable transition-function argument of
Kolmogorov~\cite{Chung67,Kendall55}. Applied to positive stability,
it also forces each monomial image under the generator to have finite
support and to change total degree by at most one.

Theorem~\ref{ana:weighted-generation} gives a necessary and sufficient
consistent-cutoff bound for the realization of the relevant formal
operators, proves uniqueness and a polynomial graph core, and constructs
compatible weighted coefficient spaces. These spaces justify the
scaling argument and first variations at arbitrary real contacts.
Theorem~\ref{ana:killed-approximation} then passes from finite capacities
to mixed capacities in the strong operator topology. Additional
generation criteria, including a complete one-variable criterion,
illustrate why preservation and bounded realization must be kept
distinct.

Section~\ref{pre:section} fixes the conventions and records the two
finite-box results imported for comparison and first-order algebra.
Section~\ref{pos:section} proves the positive classifications and their
consequences. Section~\ref{anc:section} proves the ancillary test.
Section~\ref{ana:section} gives the domain, generation, and approximation
arguments used in the infinite-capacity proofs. No spectral or quantum
algorithm result from the companion manuscript is used.


\section{Conventions and finite-box input}\label{pre:section}

Write \(\N=\{0,1,2,\ldots\}\) and
\(\HH=\{z\in\C:\Im z>0\}\). A nonzero complex polynomial is
\emph{stable} if it has no zero when all variables belong to \(\HH\).
It is \emph{real stable} when its coefficients are real, and
\emph{positive stable} when these coefficients are nonnegative.
Zero is adjoined only when a closed cone or a stable-or-zero conclusion
is explicitly used. For an operator on polynomials, preservation means
that every stable input has stable or zero output. The preserving
semigroups considered below have nonzero output on nonzero positive
inputs, as proved in the classification arguments.

For \(\kappa\in(\N_{>0}\cup\{\infty\})^d\), put
\[
 \Omega_\kappa=\{x\in\N^d:x_i\le\kappa_i\},\qquad
 X_\kappa=\ell^1(\Omega_\kappa),\qquad
 \|f\|_1=\sum_{x\in\Omega_\kappa}|f_x|.
\]
We identify a coefficient vector with its power series
\(f(z)=\sum_x f_xz^x\), initially on the unit polydisc. A vector
in \(X_\kappa\) need not define an entire function. All coefficients
are ordinary monomial coefficients. Matrix entries follow the column
convention \(Az^x=\sum_y A_{y,x}z^y\). A positive operator has
nonnegative matrix entries; a \emph{Metzler} generator matrix has
nonnegative off-diagonal entries. Markov normalization means
preservation of evaluation at \(\boldsymbol1\), or equivalently
preservation of total coefficient mass.

Capacities bound the degree separately in each coordinate. They do not
require that the upper degree be attained. Coordinates of capacity zero
are removed. When every capacity is finite, write
\(V_\kappa^{\mathbb F}=\mathbb F[z_1,\ldots,z_d]_{\deg_i\le\kappa_i}\),
with \(\mathbb F=\R\) or \(\C\). For multi-indices,
\(\partial^\alpha=\prod_i\partial_i^{\alpha_i}\) and
\((x)_\alpha=\prod_i x_i(x_i-1)\cdots(x_i-\alpha_i+1)\).
The canonical order of an operator refers to its unique normal-ordered
expression; its triangular construction is \eqref{pos:normal-recursion}.
At a binary site, \(\partial_i^2\) is absent. For a symmetric matrix
\(Q\), the convention \(\sum_{i,j}Q_{ij}\partial_i\partial_j\)
has full mixed coefficient \(2Q_{ij}\) when \(i<j\).

An inert extension acts as the identity on additional polynomial
variables. \emph{Complete complex stability preservation} means
preservation after every finite collection of inert variables with
arbitrary finite degree bounds. This notion concerns polynomial zeros.
The positive ancillary question allows only nonnegative input
coefficients and is formulated separately in Section~\ref{anc:section}.

We use real specialization, differentiation, and limits in the usual
stable-or-zero sense, and use normalized polarization
\[
 z_i^a\longmapsto \binom{\kappa_i}{a}^{-1}
 e_a(z_{i1},\ldots,z_{i\kappa_i}).
\]
Polarization preserves stability and nonnegative coefficients; diagonal
specialization recovers the original polynomial. The relevant symbol,
polarization, and closure results are cited where used
\cite{BB09,BB10,BB10II,BBL09,LVR10}.

The following two results are imported from the companion manuscript.
Their statements are included to specify all dependencies. The first
is used only to compare finite-box positive ancillary preservation with
complete complex preservation; the second is used in the mixed-capacity
growth reduction.

\begin{theorem}[Real finite-box generators, \cite{WeiGenerators2026}]
\label{gen:main-real}
Let every \(\kappa_i\) be a finite positive integer, and let
\(A\in\operatorname{End}_{\R}(V_\kappa^{\R})\). The following
conditions are equivalent:
\begin{enumerate}
\item \(e^{tA}\) preserves real stability for every \(t\ge0\).
\item Its complexification preserves complex stability for every
      \(t\ge0\), also after adjoining arbitrary inert polynomial variables.
\item The canonical representation of \(A\) has order at most two,
      \(P_{2e_i}(s)\le0\) for every \(s\in\R\) and every
      \(\kappa_i\ge2\), and \(P_{e_i+e_j}(s,t)\le0\) for
      every \((s,t)\in\R^2\) and \(i<j\).
\end{enumerate}
In the third condition box preservation forces \(P_{2e_i}\) to
depend only on \(z_i\), with degree at most four, and
\(P_{e_i+e_j}\) to depend only on \(z_i,z_j\), with bidegree
at most \((2,2)\). There are no further inequalities on the
zeroth- and first-order coefficients beyond the box-endomorphism
identities.
\end{theorem}

\begin{proposition}[First-order box structure, \cite{WeiGenerators2026}]
\label{gen:box-decomposition}
Let every \(\kappa_i\) be a finite positive integer. Every
endomorphism of \(V_\kappa^{\mathbb F}\) of canonical order at
most one has a unique expression
\[
 c\Id+\sum_{i=1}^d
 \left(q_i(z_i)\partial_i-\frac{\kappa_i}{2}q_i'(z_i)\right),
 \qquad c\in\mathbb F,\quad q_i\in\mathbb F[z_i]_{\le2}.
\]
Here \(\mathbb F\) can be \(\R\) or \(\C\).
\end{proposition}

The last statement is the first-order part of the companion proposition
``Operators of order at most two on a box.'' It is an algebraic statement,
without a positivity or semigroup assumption.


\section{Positive coefficients and particle capacities}
\label{pos:section}

The restriction to nonnegative coefficients changes the generator problem substantially.
In particular, collective particle loss is compatible with stability, whereas its extension
by an arbitrary unchanged coordinate need not be.  We first classify the unrestricted
positive problem and then give a capacity-sensitive criterion.  Throughout this section,
coefficients are ordinary monomial coefficients, and matrix entries use the convention
\(Az^x=\sum_y A_{y,x}z^y\).

Write \(\mathcal S_+\) for the nonzero real stable polynomials with nonnegative
coefficients.  Write \(\mathcal{LP}_+\) for their nonzero locally uniform entire limits.
When taking closed cones we adjoin zero explicitly.  On a countable state space we use
the coefficient space \(\ell^1\).  Thus every member of \(\mathcal{LP}_+\) belongs to
\(\ell^1\), although a general member of \(\ell^1\) need not be entire.
\begin{lemma}[Positive Laguerre--P\'olya coefficient closure, \cite{LVR10}]\label{pos:lp-closure}
Coefficientwise limits of members of \(\mathcal{LP}_+\), when the limiting
coefficient sum is finite, belong to \(\mathcal{LP}_+\cup\{0\}\).
Every member of \(\mathcal{LP}_+\) has stable Jensen-polynomial approximants
converging in \(\ell^1\), whose coefficients are bounded by those of the
limiting series.
\end{lemma}
This is the positive Laguerre--P\'olya closure and Jensen characterization.
We use this precise standard theorem when passing from coefficient norm convergence
to an entire stable limit.

\subsection{Root contacts and Poisson tangency}

\begin{lemma}[A repeated-root first variation]\label{pos:root-contact}
Suppose \(p_t\) is holomorphic near zero, has real Taylor coefficients, and
\[
p_t(u)=c u^m+t h(u)+o(t),\qquad c>0,
\]
locally uniformly as \(t\downarrow0\).  Suppose that all zeros sufficiently near zero
are real.  Then \(h(0)=0\) if \(m\ge3\), and \(h(0)\le0\) if \(m=2\).
\end{lemma}
\begin{proof}
Fix a small circle isolating the zero of \(p_0\).  The argument principle and
\(p_t'/p_t\) express the power sums of the enclosed \(m\) roots as contour integrals.
Consequently their monic root polynomial has a right first variation, and its first
two nonleading coefficients are \(O(t)\).  If its roots are \(r_1(t),\ldots,r_m(t)\),
Newton's identity gives \(\sum_jr_j(t)^2=O(t)\).  Since these roots are real, each is
\(O(\sqrt t)\).  Their product is \(o(t)\) when \(m\ge3\).  The complementary analytic
factor tends to \(c\), so this proves the first assertion.  For \(m=2\), the
discriminant of the root polynomial is nonnegative.  Its linear coefficient is
\(O(t)\), hence its constant coefficient has upper first derivative at most zero.
Multiplication by the positive limiting factor \(c\) proves the second assertion.
\end{proof}

For a capacity vector \(\kappa\in(\N\cup\{\infty\})^d\), let
\(\Omega_\kappa=\{x\in\N^d:x_i\le\kappa_i\}\).  Coordinates of capacity zero are
removed.  Every linear map on the allowed polynomials has a unique locally finite
normal-ordered expression
\begin{equation}\label{pos:normal-recursion}
 A=\sum_{\alpha\le\kappa}P_\alpha(z)\partial^\alpha,
 \qquad
 P_\alpha=\frac1{\alpha!}\left(Az^\alpha-
 \sum_{\substack{\beta\le\alpha\\\beta\ne\alpha}}
 (\alpha)_\beta z^{\alpha-\beta}P_\beta\right).
\end{equation}
Here the \(P_\alpha\) may have degrees outside the original coordinate box.
The expression is locally finite because only finitely many derivatives act on a
fixed polynomial.

\begin{lemma}[Forced differential order]\label{pos:order}
Suppose the polynomial action of a positive stability-preserving semigroup has a
locally uniform generator first variation near a real open negative box.
Then \(P_\alpha=0\) whenever \(|\alpha|\ge3\).
For unrestricted coordinates it is enough to assume preservation of every
\((c+a\cdot z)^m\), where \(c\ge0\), \(a>0\), and \(m\ge0\).
\end{lemma}
\begin{proof}
For an allowed \(\alpha\), use
\(f(z)=\prod_i(z_i+s_i)^{\alpha_i}\), with \(s_i>0\) small.
At \(z=-s\), every normal-ordered term vanishes except the term with derivative
\(\alpha\); thus \(Af(-s)=\alpha!P_\alpha(-s)\).
Along a strictly positive direction through \(-s\), the initial polynomial has a
root of multiplicity \(|\alpha|\), with positive leading coefficient.
Lemma~\ref{pos:root-contact} annihilates \(P_\alpha\) on an open box, hence
identically by analytic continuation.  For the last assertion choose \(y<0\),
\(c=-a\cdot y\), and evaluate \(A(c+a\cdot z)^m\) at \(y\).  This gives
\(m!\sum_{|\alpha|=m}a^\alpha P_\alpha(y)=0\) for every \(a>0\).
Coefficient comparison in \(a\) gives the same conclusion.
\end{proof}

The local first variation in this lemma is automatic in finite dimension.  For
an existing positive \(C_0\)-semigroup, the required domain and local analytic
statements are supplied by Theorem~\ref{ana:domain}; the subsequent extension
to arbitrary real contacts is kept separate from this initial local argument.

\begin{lemma}[Poisson tangency]\label{pos:poisson}
Let \(G_t\in\mathcal{LP}_+\), \(G_0(z)=c\exp(a\cdot z)\), where \(c>0\), \(a\ge0\).
Suppose \(G_t=G_0+tH+o(t)\) locally uniformly near the origin.  Then
\(H/G_0\) is a polynomial of total degree at most two.  Its homogeneous quadratic
part is nonpositive on every strictly positive vector.
\end{lemma}
\begin{proof}
The one-variable positive Laguerre--P\'olya factorization, including its exponential
factor, is
\[
 g(s)=g(0)e^{\sigma s}\prod_j(1+\rho_js),
 \qquad \sigma,\rho_j\ge0,\quad \sum_j\rho_j<\infty.
\]
It follows from the standard Laguerre--P\'olya product representation
\cite{LVR10}; positive coefficients put every finite zero on the negative real axis.
If \(\log(g(s)/g(0))=\sum_{k\ge1}\ell_ks^k\), then for \(k\ge2\)
\[
 \ell_k=\frac{(-1)^{k-1}}k\sum_j\rho_j^k,
 \qquad
 |\ell_k|\le\frac{(-2\ell_2)^{k/2}}k.
\]
The second inequality follows from
\(\sum\rho_j^k\le(\sum\rho_j^2)^{k/2}\).
Apply this to \(G_t(sv)\) for fixed \(v>0\).  Its logarithm has a first variation,
\(\ell_2(t)=O(t)\), and \(\ell_2(t)\le0\).  Therefore
\(\ell_k(t)=o(t)\) for every \(k\ge3\).
The first variation of its logarithm is \(H(sv)/G_0(sv)\).
Each homogeneous Taylor coefficient of degree at least three vanishes on the open
positive orthant, and is therefore the zero polynomial.  The degree-two coefficient
is nonpositive there, as claimed.
\end{proof}

\subsection{Collective loss and an exact Euler step}

Put \(E=\sum_i z_i\partial_i\), \(D=\sum_i\partial_i\), and
\begin{equation}\label{pos:coal-operator}
 C=ED-E(E-1)=\sum_{i,j}z_j(1-z_i)\partial_i\partial_j.
\end{equation}
For \(|x|=n\),
\(Cz^x=(n-1)\sum_i x_i z^{x-e_i}-n(n-1)z^x\).
Thus \(C\) is a Markov generator of single-particle losses, with the rate of losing
a type \(i\) particle equal to \(x_i(|x|-1)\).

\begin{proposition}[Sharp positive Euler step]\label{pos:coal}
On polynomials of total degree at most \(N\ge2\), the operator \(I+hC\) preserves
\(\mathcal S_+\) and evaluation at \(\boldsymbol1\) whenever
\(0\le h\le1/[N(N-1)]\).  This is the exact interval on which it is a Markov Euler
step on the whole cutoff.  Consequently \(e^{tC}\) preserves \(\mathcal S_+\).
\end{proposition}
\begin{proof}
Write \(f=\sum_{n=0}^Nf_n\) in homogeneous parts and homogenize positively:
\(F(w,z)=\sum_nw^{N-n}f_n(z)\).  Positive homogenization preserves stability
\cite[Theorem 4.5]{BBL09}.  Hence
\[
 G(\xi,\eta,z)=F(1+\xi+\eta,z+\eta\boldsymbol1)
\]
is stable.  Set \(a=1-hN(N-1)\), \(b=h(N-1)\), and
\(\mathcal L=a+b(\partial_\xi+\partial_\eta)-h\partial_\xi\partial_\eta\).
For \(h>0\) its negative differential symbol is
\(a-b(u+v)-huv\).  Solving for \(v\) gives
\[
 \Im\frac{a-bu}{b+hu}
 =-\frac{b^2+ah}{|b+hu|^2}\Im u<0,
 \qquad b^2+ah=h[1-h(N-1)]>0.
\]
Thus this symbol is stable.  The constant-coefficient differential-symbol theorem
\cite{BB09,BB10} shows that \(\mathcal LG\) is stable or zero, including in the
unchanged \(z\)-variables.
For \(k=N-n\), direct differentiation of
\((1+\xi+\eta)^kf_n(z+\eta\boldsymbol1)\) gives, at \(\xi=\eta=0\),
\[
 \mathcal LG_n=
 [a+2bk-hk(k-1)]f_n+(b-hk)Df_n
 =[1-hn(n-1)]f_n+h(n-1)Df_n.
\]
Therefore \((\mathcal LG)|_{\xi=\eta=0}=(I+hC)f\).  Real specialization preserves
stability or zero.  The displayed monomial formula for \(C\) proves nonnegativity
of the coefficients and preservation of their sum, so the output is nonzero.
For \(h>1/[N(N-1)]\), the coefficient of \(z_1^N\) in \((I+hC)z_1^N\) is negative.
Finally take Euler limits on each fixed finite total-degree cutoff; coefficient
closure and preservation of mass give the assertion for \(e^{tC}\).
\end{proof}

\subsection{The unrestricted logarithmic symbol}

For \(a\in\R^d\), define the logarithmic symbol, when the indicated exponential
belongs to the generator domain, by
\(\Psi_A(z,a)=e^{-a\cdot z}A(e^{a\cdot z})\).

\begin{theorem}[Unrestricted positive classification]\label{pos:symbol}
Let \((T_t)_{t\ge0}\) be a coefficientwise-positive \(C_0\)-semigroup of bounded
operators on \(\ell^1(\N^d)\), with generator \(A\).  The following are equivalent:
\begin{enumerate}
\item \(T_t\mathcal{LP}_+\subseteq\mathcal{LP}_+\) for every \(t\ge0\).
\item \(T_t\mathcal S_+\subseteq\mathcal{LP}_+\) for every \(t\ge0\).
\item All polynomials belong to \(D(A)\), and on polynomials
\begin{equation}\label{pos:canonical}
 \begin{aligned}
 A={}&cI+\lambda\cdot z+
 \sum_i\left(b_i+\sum_jm_{ij}z_j\right)\partial_i
 \\
 &+ED_\nu-\beta E(E-1)-\sum_i\eta_i z_i^2\partial_i^2,
 \qquad D_\nu=\sum_i\nu_i\partial_i .
 \end{aligned}
\end{equation}
where \(\lambda_i,b_i,\nu_i,\beta,\eta_i\ge0\), \(m_{ij}\ge0\) for \(i\ne j\),
and \(c,m_{ii}\in\R\).
\end{enumerate}
For \(d\ge2\) the parameters are unique.  For \(d=1\), combine \(\beta+\eta_1\)
into one nonnegative parameter.  Equivalently,
\begin{equation}\label{pos:log-symbol}
 \Psi_A(z,a)=c+\lambda\cdot z+b\cdot a+a^{\mathsf T}Mz
 +(\nu\cdot a)(z\cdot a)-\beta(z\cdot a)^2
 -\sum_i\eta_i(z_i a_i)^2.
\end{equation}
The theorem concerns an existing semigroup.  The separate existence criterion for
a formal expression is Theorem~\ref{ana:weighted-generation}.
\end{theorem}
\begin{proof}
The first implication is immediate.  Assume the second statement.  Jensen
approximation and Lemma~\ref{pos:lp-closure} give preservation of all
\(\mathcal{LP}_+\), initially allowing zero.  Zero is impossible: for each basis
vector, strong continuity gives \((T_{t/n})_{x,x}>0\) for sufficiently large
\(n\), and positivity and the semigroup law give
\((T_t)_{x,x}\ge((T_{t/n})_{x,x})^n>0\).
Every nonzero positive input dominates a positive multiple of a basis vector.
The domain conclusion follows from Corollary~\ref{ana:polynomial-domain}.  Lemma~\ref{pos:order} gives
\(A=\sum_{|\alpha|\le2}P_\alpha\partial^\alpha\), with \(P_\alpha\in\ell^1\).
Consequently
\begin{equation}\label{pos:graph-bound}
 \|Az^x\|_1\le\sum_{|\alpha|\le2}(x)_\alpha\|P_\alpha\|_1
 \le C(1+|x|)^2.
\end{equation}
The polynomial partial sums of \(e^{a\cdot z}\) converge in graph norm by this
bound, since Poisson factorial moments are finite.  Closedness of \(A\) gives
\[
 e^{a\cdot z}\in D(A),\qquad
 A(e^{a\cdot z})=e^{a\cdot z}\sum_{|\alpha|\le2}a^\alpha P_\alpha(z).
\]
The stable approximants \((1+a\cdot z/m)^m\) converge in coefficient norm.
Their images and positive coefficient closure imply
\(T_te^{a\cdot z}\in\mathcal{LP}_+\).
Lemma~\ref{pos:poisson}, followed by coefficient comparison in \(a>0\), now gives
\(\deg P_\alpha\le2\) and
\begin{equation}\label{pos:poisson-inequality}
 P_0^{[2]}(v)+\sum_i a_iP_i^{[2]}(v)
 +\sum_{|\alpha|=2}a^\alpha P_\alpha^{[2]}(v)\le0
 \quad(a,v>0).
\end{equation}

Positivity of the semigroup makes the off-diagonal entries of \(A\) nonnegative.
The nonconstant coefficients of \(P_0=A1\) are therefore nonnegative.  Let
\(a\to0\) in \eqref{pos:poisson-inequality}; its quadratic part must vanish, giving
\(P_0=c+\lambda\cdot z\), \(\lambda\ge0\).
If \(|\rho|=2\), the coefficient of \(z^{(n-1)e_i+\rho}\) in \(Az_i^n\) is
\(n[z^\rho]P_i+O(1)\).  No second-order term has sufficiently high output degree
to contribute.  Positivity for arbitrarily large \(n\) gives
\([z^\rho]P_i\ge0\).  Replacing \(a\) by \(\varepsilon a\) in
\eqref{pos:poisson-inequality} and taking \(\varepsilon\downarrow0\) eliminates
these quadratic coefficients.  Thus
\(P_i=b_i+\sum_jm_{ij}z_j\), with the asserted signs obtained from \(Az_i\).

At this point \(A-M_\lambda\), where \(M_\lambda f=(\lambda\cdot z)f\), preserves
every finite total-degree cutoff, and \(M_\lambda\) is bounded on \(\ell^1\).
Theorem~\ref{ana:weighted-generation} gives compatible weighted semigroups:
if the cutoff-preserving semigroup has bound \(Me^{\omega t}\), then
\(\|T_tf\|_R\le M R^N e^{(\omega+M\Lambda R)t}\|f\|_1\) for
\(\deg f\le N\), \(R\ge1\), where \(\Lambda=\sum_i\lambda_i\).
Since \(A^2f\) has degree at most \(N+2\), the Taylor integral gives
\[
 \|T_tf-f-tAf\|_R\le
 \frac{t^2}{2}M R^{N+2}e^{(\omega+M\Lambda R)\tau}\|A^2f\|_1
 \quad(0\le t\le\tau).
\]
Thus generator first variations are legitimate at every real point, not only
in the unit polydisc.  We now determine the second-order part algebraically.

For \(|\alpha|=2\),
\begin{equation}\label{pos:reaction-positive}
 [z^\rho]P_\alpha\ge0\qquad(\rho\ne\alpha).
\end{equation}
For \(\alpha=2e_i\), divide the relevant off-diagonal coefficient in \(Az_i^n\)
by \(n(n-1)\) and let \(n\to\infty\).  For \(\alpha=e_i+e_j\), at least one of
\(\rho_i,\rho_j\) is zero if \(|\rho|\le2\) and \(\rho\ne\alpha\).
If \(\rho_i=0\), use source \(e_i+ne_j\) and target \((n-1)e_j+\rho\).
Terms differentiating only \(z_j\) retain \(z_i\), and \(\partial_i^2\) vanishes.
The coefficient is therefore \(n[z^\rho]P_\alpha+O(1)\).  This proves
\eqref{pos:reaction-positive}.

The square test \((c+a\cdot z)^2\) gives
\[
 H_a(z):=\sum_{|\alpha|=2}a^\alpha P_\alpha(z)\le0
 \quad\text{whenever }a>0,\ a\cdot z<0.
\]
Scaling \(z\) to zero first removes the nonnegative constant terms.
The linear part \(\ell(a)\cdot z\) is nonpositive on the half-space
\(a\cdot z<0\), so \(\ell(a)=k(a)a\) with \(k(a)\ge0\).
The polynomial identities \(a_j\ell_i=a_i\ell_j\) show that every \(\ell_i\)
is divisible by \(a_i\), and the linear quotients agree.  Thus
\(k(a)=\nu\cdot a\), with \(\nu_i\ge0\).
Scaling \(z\) to infinity shows that the quadratic part of \(H_a\) is nonpositive
on the same half-space.  Being even in \(z\), it is nonpositive everywhere.

Here is the full elimination of off-diagonal quadratic outputs.  Set \(a=e_i\).
On \(z_i=0\), the nonnegative coefficients in \eqref{pos:reaction-positive}
force every quadratic term not containing \(z_i\) to vanish.  A surviving
\(z_i z_j\) term is also impossible, because \(z_j\) can have either sign and
arbitrarily large magnitude.  For mixed sources, the coefficient of \(z_k^2\)
in the nonpositive quadratic form is
\(-h_{kk}a_k^2+\sum_{i<j}[z_k^2]P_{e_i+e_j}a_i a_j\).
Put \(a_k=\varepsilon\), all other entries equal to one, and let
\(\varepsilon\downarrow0\).  This first removes all coefficients whose
source avoids \(k\).  Divide the remaining inequality by \(\varepsilon\)
and take the limit again to remove sources containing \(k\).  Thus every
displayed nonnegative mixed-source coefficient vanishes.  Finally, a mixed output different from its mixed source
has an index outside the source.  Support \(a\) on the two source indices.
The corresponding quadratic form has zero square coefficient at the outside
index; a nonpositive quadratic form with such a zero has an identically zero
row there.  This removes that output as well.  We have proved
\[
 P_{2e_i}=\nu_i z_i-h_{ii}z_i^2,\qquad
 P_{e_i+e_j}=\nu_jz_i+\nu_i z_j-2h_{ij}z_i z_j,
\]
where \(H=(h_{ij})\) is positive semidefinite.

Write the principal matrix as
\(Q(z)=(\nu z^{\mathsf T}+z\nu^{\mathsf T})/2-
\diag(z)H\diag(z)\).  Different affine factors supply the additional restriction.
At a real common zero of
\((c_1+a\cdot z)(c_2+b\cdot z)\), with nonnegative factors,
Lemma~\ref{pos:root-contact} gives \(a^{\mathsf T}Q(z)b\le0\).
For distinct \(i,j,k\), apply it to \((z_i+z_k)(1+z_j)\) at
\((z_i,z_k,z_j)=(s,-s,-1)\).  The result is
\[
 -\frac{\nu_i+\nu_k}{2}+s(h_{ij}-h_{kj})\le0\quad(s\in\R).
\]
Hence all off-diagonal entries of \(H\) are a common value \(\beta\) when
\(d\ge3\); for \(d=2\), put \(\beta=h_{12}\).
The contact \((s+z_i)(s+z_j)\), at \(z_i=z_j=-s\), gives
\(-s(\nu_i+\nu_j)/2-\beta s^2\le0\), whence \(\beta\ge0\).
The contact \((z_i+z_j)(s+z_i)\), at \((z_i,z_j)=(-s,s)\), gives
\(-s(\nu_i+\nu_j)/2-s^2(h_{ii}-\beta)\le0\), whence \(h_{ii}\ge\beta\).
Thus \(H=\beta\boldsymbol1\boldsymbol1^{\mathsf T}+\diag(\eta)\),
\(\beta,\eta\ge0\).  This proves necessity and uniqueness.

For sufficiency, remove the bounded multiplication \(M_\lambda\) and work first
on a fixed total-degree cutoff.  The affine first-order part has flow
\[
 f(z)\longmapsto e^{ct}
 f\left(e^{tM}z+\int_0^te^{sM}b\,ds\right).
\]
The matrix \(M\) is Metzler, so \(e^{tM}\) is nonnegative with positive diagonal;
this substitution preserves \(\mathcal S_+\).
For \(\beta>0\) and \(\nu_i>0\), conjugate \(\beta C\) by positive coordinate
scaling \(S_sf(z)=f(s_1z_1,\ldots,s_dz_d)\), with \(s_i=\nu_i/\beta\):
\[
 S_s^{-1}(\beta C)S_s=ED_\nu-\beta E(E-1).
\]
Proposition~\ref{pos:coal} proves preservation.  Zero parameters follow by
parameter limits on the same finite-dimensional cutoff.
For individual damping \(K_i=-z_i^2\partial_i^2\), the active algebraic symbol of
\(I+hK_i\) at coordinate degree \(N\) factors as
\[
 (z+w)^{N-2}
 [w+(1-\sqrt\theta)z][w+(1+\sqrt\theta)z],
 \qquad\theta=hN(N-1)\in[0,1].
\]
It is stable, and the coefficients of the operator on monomials are nonnegative.
The finite-degree symbol theorem and Euler limits give preservation by
\(e^{t\eta_iK_i}\).  Finite-dimensional Lie products combine these elementary
flows.  Stability follows by closure, and invertibility of each finite
exponential excludes zero outputs.
Finally the bounded-perturbation product formula, justified in
Theorem~\ref{ana:weighted-generation}, adds \(M_\lambda\); its own flow is
multiplication by \(e^{t\lambda\cdot z}\).  Jensen approximation and coefficient
closure extend the result to \(\mathcal{LP}_+\).  Nonzero outputs follow from
\(T_tz^x\ge e^{A_{x,x}t}z^x\), the variation-of-constants inequality for a positive
semigroup and a basis vector in its domain.
\end{proof}

\begin{corollary}[Arbitrary-jump Markov rigidity]\label{pos:markov}
Let \(T_t\) be the transition semigroup of a conservative nonexplosive continuous-time
Markov chain on \(\N^d\), with finite exit rate at every state.  No restriction on
jump sizes or rate dependence is imposed.  The following are equivalent:
\begin{enumerate}
\item Every input in \(\mathcal S_+\) is mapped into \(\mathcal{LP}_+\) at every time.
\item The same assertion holds for all \((c+a\cdot z)^m\), where
\(c\ge0\), \(a>0\), and \(m\ge0\).
\item The generator is
\begin{equation}\label{pos:markov-generator}
 A=\sum_i\lambda_i(z_i-1)+
 \sum_i\left[\delta_i(1-z_i)+\sum_{j\ne i}r_{ij}(z_j-z_i)\right]\partial_i
 +\gamma C,
\end{equation}
with \(\lambda_i,\delta_i,r_{ij},\gamma\ge0\).
\end{enumerate}
In particular, the only transitions are constant immigration, independent
migration, and the single deaths
\(x\to x-e_i\) at rate \(x_i[\delta_i+\gamma(|x|-1)]\).
If all outputs from polynomial inputs must remain polynomials, then
\(\lambda=0\), and this condition is sufficient.
\end{corollary}
\begin{proof}
Finite exit rates put every basis vector in the \(\ell^1\) generator domain:
conditioning on the first jump gives
\(\|(T_tz^x-z^x)/t-Az^x\|_1\to0\) and
\(\|T_tz^x-z^x\|_1\le2tq(x)\).
Under the second assertion, the multinomial version of Lemma~\ref{pos:order}
applies.  The diagonal coefficient gives
\[
 q(x)=-\sum_{|\alpha|\le2}(x)_\alpha[z^\alpha]P_\alpha
 \le C(1+|x|)^2.
\]
Normalized multinomial approximants to a product-Poisson law converge in total
variation.  Markov contraction and positive coefficient closure therefore imply
preservation of those Poisson inputs.  The graph-norm and Poisson-tangency argument
in Theorem~\ref{pos:symbol} applies unchanged.  It yields constant immigration,
affine first-order coefficients, and a second-order part
\[
 P_{2e_i}=\nu_i z_i-h_{ii}z_i^2,\qquad
 P_{e_i+e_j}=\nu_jz_i+\nu_i z_j-2h_{ij}z_i z_j,
 \qquad H\succeq0.
\]
Only square multinomial contacts were used to reach this point.
Mass preservation gives \(P_\alpha(\boldsymbol1)=0\), hence
\(h_{ii}=\nu_i\) and \(2h_{ij}=\nu_i+\nu_j\).
The quadratic part of the square-contact expression is consequently
\(-(a\cdot z)(\sum_i\nu_i a_i z_i)\), nonpositive for every real \(z\).
Two nonzero real linear forms whose product is everywhere nonnegative are
nonnegative multiples of one another: the second must vanish on the kernel
of the first.  Taking \(a=\boldsymbol1\) shows that every \(\nu_i\) equals a
common \(\gamma\ge0\).  The constant and linear mass identities give exactly
\eqref{pos:markov-generator}.  This proves the difficult implication without
assuming preservation of products of different affine factors.
Sufficiency follows from Theorem~\ref{pos:symbol}, or from the same elementary
flows and bounded immigration construction.  The chain in
\eqref{pos:markov-generator} is nonexplosive: its population is bounded by the
initial population plus a rate \(\sum_i\lambda_i\) Poisson immigration process,
and every bounded-population set is finite.
If \(\lambda_i>0\), from the empty state a path with any prescribed number of
type \(i\) immigrations and no other jumps has positive probability before any
fixed positive time.  The output therefore has infinite support.  Conversely
\(\lambda=0\) leaves every total-degree cutoff invariant.
\end{proof}

Modulo the lineality generated by \(I,z_1\partial_1,\ldots,z_d\partial_d\),
Theorem~\ref{pos:symbol} describes, for \(d\ge2\), the formal cone generated by
\[
 z_i,\quad\partial_i,\quad z_j\partial_i\ (i\ne j),\quad
 E\partial_i,\quad-E(E-1),\quad-z_i^2\partial_i^2.
\]
This is a statement about the formal cone; it does not assert that each ray
individually generates bounded operators on every fixed coefficient space.

\begin{corollary}[Diagonal selection]\label{pos:diagonal}
Suppose \(v:\N^d\to\R\) is bounded above.  The positive semigroup
\(T_tz^x=e^{tv(x)}z^x\) preserves \(\mathcal{LP}_+\) if and only if
\[
 v(x)=c+\mu\cdot x-\beta|x|(|x|-1)-\sum_i\eta_i x_i(x_i-1),
 \qquad \beta,\eta_i\ge0,
\]
with the same one-dimensional merging convention as above.
\end{corollary}
\begin{proof}
The upper bound gives a bounded-operator positive \(C_0\)-semigroup on \(\ell^1\).
In \eqref{pos:canonical}, every nondiagonal parameter must vanish because its
monomial matrix entry is nonnegative and cannot cancel another off-diagonal
entry.  The remaining diagonal terms are exactly the displayed expression.
The converse is Theorem~\ref{pos:symbol}.
\end{proof}

\begin{example}[The square-symbol test is insufficient]\label{pos:psd-counterexample}
The diagonal contraction generated by
\(A=-z_1^2\partial_1^2-4z_1z_2\partial_1\partial_2-4z_2^2\partial_2^2\)
has logarithmic symbol \(-(a_1z_1+2a_2z_2)^2\).
Nevertheless, at \(t=(\log2)/2\), it sends the stable polynomial
\((z_1+z_2)(1+z_1)\) to
\(z_1+z_2+z_1^2/2+z_1z_2/4\), which vanishes at
\((z_1,z_2)=(-3+i,2+6i)\).
Thus positive semidefiniteness of the preliminary damping matrix is strictly
weaker than the common-damping condition in Theorem~\ref{pos:symbol}.
By contrast, \(E(\partial_1+2\partial_2)-2E(E-1)\) is a positive contraction
preserving stability: its column sum at \(z^x\) is
\(-x_1(|x|-1)\le0\), and it has the form \eqref{pos:canonical}.
It represents unequal collective loss together with killing.
\end{example}

\subsection{Finite capacities and quadratic contacts}

Now let every \(\kappa_i\) be finite and positive, put
\(V_\kappa=\operatorname{span}\{z^x:x\le\kappa\}\), and let
\(J=\{i:\kappa_i=1\}\).  The numerical capacities enter the boundary factors of
the rates.  The intrinsic contact domain remembers exactly the binary coordinates.
With \(e=\boldsymbol1\), define \(\mathfrak Q_J\) by
\begin{equation}\label{pos:contact-domain}
 \begin{gathered}
 y\in\R^d,\quad H=H^{\mathsf T},\quad H_{ij}\ge0,
 \quad H_{ii}=0\ (i\in J),\quad e^{\mathsf T}He=1,\\
 (He)(He)^{\mathsf T}-H\succeq0,\qquad
 Hy\le0,\qquad y^{\mathsf T}Hy\ge0.
 \end{gathered}
\end{equation}
Here \(H_{ij}\ge0\) is entrywise nonnegativity, and is not a positive-semidefinite
condition on \(H\).  Under the normalization, the displayed matrix inequality
is equivalent to \(H\) having exactly one positive eigenvalue.  Indeed, decompose
\(u=(e^{\mathsf T}Hu)e+v\), with \(e^{\mathsf T}Hv=0\).  A quadratic form with
one positive direction is nonpositive on this orthogonal complement, which gives
\(u^{\mathsf T}Hu\le(e^{\mathsf T}Hu)^2\); the converse follows by restriction.

\begin{theorem}[Finite-capacity classification]\label{pos:finite}
A real operator \(A:V_\kappa\to V_\kappa\) generates a positive semigroup preserving
\(\mathcal S_+\cap V_\kappa\) if and only if its only off-diagonal entries are
\begin{equation}\label{pos:finite-rates}
\begin{aligned}
 A_{x-e_i,x}&=x_i\left[\delta_i+\alpha_i(x_i-1)
                         +\sum_{j\ne i}a_{ij}x_j\right],\\
 A_{x+e_i,x}&=(\kappa_i-x_i)\left[b_i+\chi_i(\kappa_i-x_i-1)
                         +\sum_{j\ne i}g_{ij}(\kappa_j-x_j)\right],\\
 A_{x-e_i+e_j,x}&=c_{ij}x_i(\kappa_j-x_j)\qquad(i\ne j),
\end{aligned}
\end{equation}
all the displayed parameters are nonnegative, and the diagonal is
\begin{equation}\label{pos:diagonal-rate}
 A_{x,x}=\omega+\sum_i\mu_ix_i+\sum_i\tau_ix_i(x_i-1)
                      +\sum_{i<j}\theta_{ij}x_ix_j,
\end{equation}
with real coefficients.  Set \(\alpha_i=\chi_i=\tau_i=0\) for \(i\in J\).
Writing the second-order part as \(\sum_{i,j}Q_{ij}(z)\partial_i\partial_j\),
where \(Q\) is symmetric, its entries are
\begin{equation}\label{pos:principal-matrix}
\begin{aligned}
 Q_{ii}(z)&=\alpha_i z_i+\tau_i z_i^2+\chi_i z_i^3,\\
 2Q_{ij}(z)&=a_{ij}z_j+a_{ji}z_i+\theta_{ij}z_iz_j
 -c_{ij}z_j^2-c_{ji}z_i^2
 +g_{ij}z_i^2z_j+g_{ji}z_iz_j^2.
\end{aligned}
\end{equation}
The additional necessary and sufficient condition is
\begin{equation}\label{pos:contact-test}
 \operatorname{tr}(Q(y)H)\le0\qquad((y,H)\in\mathfrak Q_J).
\end{equation}
Within the structural family \eqref{pos:finite-rates}--\eqref{pos:diagonal-rate},
every failure of preservation has a quadratic stable input witnessing failure
at first order in time.
\end{theorem}

We give the structural and geometric parts of the proof separately.

\begin{lemma}[Boundary factors and finite rate reduction]\label{pos:finite-structure}
Every positive stability-preserving generator on \(V_\kappa\) has the rate form
\eqref{pos:finite-rates}--\eqref{pos:diagonal-rate}, with the stated signs.
\end{lemma}
\begin{proof}
Lemma~\ref{pos:order} gives order at most two.  For a net shift \(\sigma\), write
\(p_\sigma(x)=[z^{x+\sigma}]Az^x\).  Normal ordering gives a polynomial of total
degree at most two in \(x\), represented in reduced degree on every finite
coordinate grid.  Put \(r_i=(-\sigma_i)_+\), \(s_i=(\sigma_i)_+\).  The rate
vanishes when a source has fewer than \(r_i\) particles or fewer than \(s_i\)
holes.  Reduced interpolation on each boundary slice makes these vanishings
polynomial identities.  Successive division gives
\begin{equation}\label{pos:boundary-factors}
 p_\sigma(x)=\prod_i(x_i)_{r_i}(\kappa_i-x_i)_{s_i}\,h_\sigma(x),
 \qquad\deg h_\sigma\le2-\sum_i(r_i+s_i).
\end{equation}
A shift with no admissible source has zero reduced polynomial.  Thus a nonzero
shift has \(\|\sigma\|_1\le2\).

To exclude double births, set \(g_t(s)=(e^{tA}1)(s,\ldots,s)\).
It has nonnegative coefficients, negative real zeros, and \(g_0=1\).
The coefficient of \(s^2\) in \(\log g_t\) is nonpositive.  Its derivative at
zero is \([s^2](A1)(s,\ldots,s)\), the sum of the nonnegative double-birth
rates at the empty state.  Each is zero.  Formula~\eqref{pos:boundary-factors}
shows that a double-birth rate is a constant times its two hole factors, so
it vanishes everywhere.  The involution
\(\mathcal Jf(z)=z^\kappa f(1/z)\) preserves nonnegative coefficients and real
stability: reciprocal upper-half-plane variables lie in the lower half-plane,
where a real stable polynomial is also nonvanishing.  Conjugating by \(\mathcal J\)
therefore excludes double deaths.  Only single births, single deaths, transfers,
and a quadratic diagonal remain.

The single-death affine factor has a unique expansion in
\(1,x_i-1,x_j\ (j\ne i)\), and the birth factor has a unique expansion in the
corresponding hole coordinates; ineffective binary terms are set to zero.
The transfer boundary factors exhaust the degree budget, giving \(c_{ij}\ge0\).
The pair test \((z_i+s)(z_j+t)\) at \((-s,-t)\) gives
\[
 c_{ji}s^2+c_{ij}t^2+a_{ji}s+a_{ij}t-\theta_{ij}st
 +g_{ij}s^2t+g_{ji}st^2\ge0\qquad(s,t>0).
\]
Let \(s\downarrow0\) and then \(t\downarrow0\) to obtain \(a_{ij}\ge0\);
exchange indices for \(a_{ji}\).  Let \(s\to\infty\) with \(t\) fixed, then
\(t\to\infty\), to obtain \(g_{ij}\ge0\); exchange indices again.
For \(\kappa_i\ge2\), the square test gives
\(\alpha_i-\tau_i s+\chi_i s^2\ge0\) for \(s>0\), so
\(\alpha_i,\chi_i\ge0\).
A state with one particle at site \(i\) and no others gives \(\delta_i\ge0\).
A state with exactly one hole at \(i\) and no other holes gives \(b_i\ge0\).
These are all asserted signs.  Direct normal ordering of the rates yields
\eqref{pos:principal-matrix}.
\end{proof}

\begin{lemma}[Quadratic witnesses]\label{pos:quadratic-witness}
For \((y,H)\in\mathfrak Q_J\), the polynomial
\(f_{y,H}(z)=(z-y)^{\mathsf T}H(z-y)\) belongs to
\(\mathcal S_+\cap V_\kappa\), and
\(Af_{y,H}(y)=2\operatorname{tr}(Q(y)H)\).
\end{lemma}
\begin{proof}
Nonnegative quadratic, linear, and constant coefficients are respectively the
conditions \(H\ge0\) entrywise, \(-2Hy\ge0\), and \(y^{\mathsf T}Hy\ge0\).
The binary restrictions are exactly \(H_{ii}=0\).
For \(v>0\), \(v^{\mathsf T}Hv>0\).  Since \(H\) has one positive eigenvalue,
its form is nonpositive on the \(H\)-orthogonal complement of \(v\).
If \((u+iv)^{\mathsf T}H(u+iv)=0\), the imaginary part makes \(u\) belong
to this complement, whereas the real part requires
\(u^{\mathsf T}Hu=v^{\mathsf T}Hv>0\), a contradiction.
Thus \(z^{\mathsf T}Hz\), and its real translate, are stable.
At \(y\) both the value and the gradient vanish, leaving precisely the asserted
second-order trace.  Lemma~\ref{pos:root-contact} proves necessity of
\eqref{pos:contact-test}; a strict violation proves the witness assertion.
\end{proof}

For sufficiency we need a precise finite-dimensional boundary argument.  We first
record explicitly the homogeneous phase property used in its localizations.

\begin{lemma}[Homogeneous common phase]\label{pos:phase}
A nonzero homogeneous complex stable polynomial is a nonzero complex scalar times
a polynomial with nonnegative real coefficients.
\end{lemma}
\begin{proof}
Let its degree be \(m\), and fix \(e>0\).  Homogeneity gives
\(P(ie)=i^mP(e)\ne0\).  For every real \(x\), the degree-\(m\) polynomial
\(P(x+te)\) has no upper-half-plane zero by stability and no lower-half-plane zero,
since \(P(x+te)=(-1)^mP(-x-te)\).  Dividing by its leading coefficient \(P(e)\)
gives a monic real-rooted polynomial.  In particular \(P(x)/P(e)\) is real for
every real \(x\), so \(P_0=P/P(e)\) has real coefficients.  It does not vanish
on the positive orthant, and equals one at \(e\), hence is positive there.
Fix positive real values of all coordinates except one.  The resulting
univariate polynomial is real-rooted, positive on the positive half-line,
and has positive leading coefficient.  Its roots are nonpositive, so all its
coefficients are nonnegative.  Each of these coefficient polynomials in the
remaining variables is homogeneous and stable or zero, by differentiation and
real specialization, and is nonnegative on the remaining positive orthant.
Induction on the number of variables therefore makes every coefficient of
\(P_0\) nonnegative.
\end{proof}

\begin{lemma}[Projective boundary and interior density]\label{pos:boundary}
Let \(K=(\mathcal S_+\cap V_\kappa)\cup\{0\}\), and let \(O\) be the open set
of polynomials with every allowed coefficient strictly positive.
A stable section of multidegree \(\kappa\) on \((\mathbb{CP}^1)^d\) having no
real projective singularity is interior to the real-stable set.
Moreover \(K=\overline{\operatorname{int}K\cap O}\).
\end{lemma}
\begin{proof}
Use real M\"obius charts of positive determinant at infinity; they preserve the
upper half-plane.  Suppose a stable section has a zero with a nonempty set \(I\)
of real coordinates and a nonempty complementary set \(J\) strictly in the upper
half-plane.  Real specialization and Hurwitz's theorem imply
\(f(x_I,z_J)\equiv0\).
Let \(h(u_I,z_J)\) be the first nonzero homogeneous term of
\(f(x_I+u_I,z_J)\) in \(u_I\).  Positive rescaling and coefficient limits show
that \(h\) is jointly stable.  For each \(z_J\in\HH^J\), its homogeneous
coefficients in \(u_I\) have one common phase.  Every nonzero such coefficient
is itself stable, by derivative and specialization closure, and is nonvanishing
on \(\HH^J\).  Ratios of any two of them are therefore holomorphic and real-valued
on this open connected set, hence constant.  Polynomial continuation gives
\(h(u_I,z_J)=h_0(u_I)g(z_J)\), where \(h_0\) is real homogeneous up to a phase
and \(g\) is a real stable section up to the opposite phase.
If the order in \(u_I\) is at least two, every point of this real slice is a
singularity.  If it is one, the first derivatives along the slice are multiples
of \(g\).  This section has a real projective zero: if it depends genuinely on
some coordinate, specialize the others generically to real values and use a
nonconstant real-rooted polynomial; if a prescribed positive degree is missing,
use infinity in that coordinate.  At such a zero all first derivatives of \(f\)
vanish.  Thus a partially nonreal boundary zero forces a real singularity.

At an all-real regular zero, linear localization shows that the nonzero gradient
components have the same sign.  If the components in a nonempty set \(J\) vanish,
the slice in those coordinates must vanish identically.  To justify this last
assertion, otherwise choose a positive direction \(v_J\) on which the slice is
not identically zero and a coordinate \(i\) with nonzero partial derivative.
The real specialization
\(q(s,t)=f(p+s e_i+t v_J)\) is stable,
\(q_s(0,0)\ne0\), and \(q(0,t)=ct^m+O(t^{m+1})\), with \(m\ge2\), \(c\ne0\).
The analytic implicit function theorem gives a zero branch
\(s=\phi(t)=-(c/q_s(0,0))t^m+O(t^{m+1})\).
For an appropriate ray \(t=\rho e^{i\vartheta}\), \(0<\vartheta<\pi\),
its leading term has positive imaginary part.  For small \(\rho>0\), both
\(s\) and \(t\) lie in the upper half-plane, contradicting stability.
The slice therefore vanishes identically and the preceding slice argument
again supplies a real projective singularity.

In the absence of singularities, every real zero thus has all gradient components
nonzero and of the same sign.  Choose that sign \(\sigma\) in a small complex
coordinate neighborhood, so that \(\Re(\sigma\partial_i f)>0\) there.
The inequality persists for sufficiently small real coefficient perturbations \(g\).
For real \(u\) and positive \(v\) within this neighborhood,
\[
 \Im(\sigma g(u+iv))=
 \sum_i v_i\int_0^1\Re(\sigma\partial_i g(u+itv))\,dt>0.
\]
Outside these neighborhoods the section is nonzero on the product of closed upper
hemispheres, by the exclusion of partially nonreal zeros.  Compactness gives a
finite cover and a positive lower bound on the complementary compact set.
Every sufficiently small real perturbation is therefore stable.  This proves
the first assertion.

To prove density, apply small independent birth--death flows, with strictly positive
birth and death rates, to each finite site.  Their two-state M\"obius matrices have
positive entries and positive determinant, so the induced degree-\(\kappa_i\)
action preserves stability.  Their transition matrices are strictly positive at
positive times, since every state communicates through nearest-neighbor moves.
Thus every nonzero member of \(K\) is approximated by members of \(K\cap O\).
For such an \(f\), apply \(D_s=e^{-sE(E-1)}\), which preserves positive stability
by Theorem~\ref{pos:symbol}'s finite-cutoff proof and leaves the box invariant.
For any real \(x\), the radial polynomial \(r\mapsto f(rx)\) is real-rooted.
Indeed, positive homogenization \(F(w,z)=w^Nf(z/w)\) is stable, and
\(t\mapsto F((0,x)+t(1,\varepsilon,\ldots,\varepsilon))\) is real-rooted for
\(\varepsilon>0\).  Its limit \(F(t,x)\) has nonzero leading coefficient
\(F(1,0)=f(0)>0\); reciprocity gives the assertion for \(f(rx)\).

Put \(p_s(r)=(D_sf)(rx)\).  It remains real-rooted and satisfies
\(\partial_s p_s=-r^2p_s''\).
Every nonzero root is simple when \(s>0\).  Otherwise let \(r_0\ne0\) have
multiplicity \(m\ge2\) at an interior time \(s_0>0\), and factor the nearby root
cluster as
\(u^m+a_1(s)u^{m-1}+a_2(s)u^{m-2}+\cdots\), \(u=r-r_0\).
Contour power sums make its coefficients differentiable.  Dividing the evolution
identity by the nonvanishing complementary factor shows that the derivative
of this monic factor has no terms below \(u^{m-2}\), and that
\(a_2'(s_0)=-r_0^2m(m-1)\).
The sum of squares of the real root displacements is \(a_1^2-2a_2\), nonnegative
and zero at \(s_0\), but with derivative \(2r_0^2m(m-1)>0\), a contradiction.
A finite critical zero of \(D_sf\) would give a repeated radial root at \(r=1\),
so none exists.  At a projective face with coordinates \(I\) at infinity, let
\(m=\sum_{i\in I}\kappa_i\) and let \(g\) be the leading-coefficient face of \(f\).
The corresponding face of \(D_sf\) is
\[
 e^{-sm(m-1)}(D_sg)(e^{-2sm}z_{I^c}).
\]
The face \(g\) is stable, has all coefficients positive and has positive constant
coefficient; the same radial argument excludes a critical zero there.  A zero
with all coordinates at infinity is excluded by the positive top coefficient.
The first assertion makes \(D_sf\) an interior point.  Letting the two smoothing
times tend to zero, and finally scaling a positive interior polynomial to zero,
proves the density assertion.
\end{proof}

\begin{lemma}[Ordinary boundary contacts]\label{pos:ordinary}
Let \(M=\dim V_\kappa\).  Except for a semialgebraic subset of dimension at most
\(M-2\), every point of \(\partial K\cap O\) has a unique real projective
singularity.  It is finite, has nonsingular Hessian, and admits a local defining
function
\[
 h(g)=\sigma g(y(g)),\qquad K\text{ locally }=\{h\le0\},
\]
where \(y(g)\) is its continued critical point and \(\sigma\in\{1,-1\}\).
At the original point \(f\), the matrix
\(H=\sigma\operatorname{Hess}f(y)/2\), after positive normalization, belongs
to \(\mathfrak Q_J\).
\end{lemma}
\begin{proof}
The univariate capacity-one case is immediate: the strictly positive affine
polynomials are interior.  In all other cases, for fixed projective \(p\), the
conditions \(f(p)=0\), \(df(p)=0\) impose \(d+1\) independent linear conditions.
First jets are surjective since every capacity is positive.  The incidence space
therefore has dimension \(M-1\).
Hessian degeneracy imposes a further nontrivial polynomial equation in each fiber:
the allowed local quadratic \(\sum_{i<j}u_iu_j\) has invertible Hessian for
\(d\ge2\); in one variable use \(u^2\).
A projective coordinate at infinity likewise reduces the dimension of the point
parameter by at least one.  At a finite nondegenerate singularity, the implicit
function theorem continues the critical point as \(y(g)\); the differential of
its critical value is evaluation at \(y\).  Evaluations at two distinct finite
points are linearly independent, as is seen on \(1,z_1,\ldots,z_d\).
Two distinct ordinary singularities thus lie on transverse critical-value
hypersurfaces and impose codimension at least two.
All these sets are projections of real algebraic incidence sets in finitely many
projective charts, hence semialgebraic.  Their dimension bounds give finite
smooth stratifications, each covered by countably many charts of dimension at
most \(M-2\).  This proves the exceptional-set assertion.

At a remaining boundary point, Lemma~\ref{pos:boundary} supplies its singularity.
The quadratic localization of the stable positive homogenization at \((1,y)\)
is, with its common sign chosen,
\((v-sy)^{\mathsf T}H(v-sy)\).
The homogeneous common-phase property and quadratic stability yield
\(H\ge0\) entrywise, \(Hy\le0\), \(y^{\mathsf T}Hy\ge0\), and exactly one
positive eigenvalue.  The box imposes \(H_{ii}=0\) for binary coordinates.
Its positive normalization is therefore in \(\mathfrak Q_J\).

The Hessian remains invertible nearby, and compactness excludes other projective
singularities outside the chosen neighborhood.  Lemma~\ref{pos:boundary} implies
that the local boundary is contained in \(h(g)=\sigma g(y(g))=0\).
Here \(dh\ne0\), since its differential is signed evaluation at \(y\).
If \(h(g)>0\), a positive-direction restriction through \(y(g)\) has a positive
critical value and a positive quadratic leading coefficient near the former
double root; its two nearby zeros are nonreal.  That side is unstable.
Interior density from Lemma~\ref{pos:boundary} gives stable interior points on
the side \(h<0\).  Shrink to an implicit-function chart in which this side is
connected.  Membership in \(K\) cannot change within the side without another
boundary point, and every boundary point has \(h=0\).
Thus the whole side \(h<0\) is stable.  Closedness then includes \(h=0\), giving
exactly the stated local half-space description.
\end{proof}

\begin{lemma}[Integration through ordinary boundary points]\label{pos:integration}
Let \(K\subset\R^M\) be closed, \(O\) open, and
\(K=\overline{\operatorname{int}K\cap O}\).
Suppose a complete ambient flow generated by a \(C^1\) vector field leaves
\(O\) forward invariant.
Suppose the boundary in \(O\), except for countably many \(C^1\) images of
dimension at most \(M-2\), has \(C^2\) local defining functions \(h\) with
\(dh\ne0\) and \(K=\{h\le0\}\).
Then the flow preserves \(K\) if and only if \(dh(V)\le0\) at every such
ordinary boundary point, where \(V\) is its vector field.
\end{lemma}
\begin{proof}
Necessity is differentiation from a boundary point.
For sufficiency fix a time horizon \(T\).  If \(\Sigma\) is the exceptional set,
the initial points whose trajectories meet \(\Sigma\) before time \(T\) lie in
the union of the images
\((t,u)\mapsto\Phi_{-t}(\psi(u))\), \(0\le t\le T\), where \(\psi\) ranges
over its charts.  Each parameter space has dimension at most \(M-1\), and
countable compact covers make these images Lipschitz images of lower-dimensional
parameter sets relative to \(M\)-dimensional volume.  Their union has measure zero.
Only the ambient inverse flow is used; \(O\) need not be backward invariant.

In a regular chart, the ordinary boundary points are dense in its smooth
hypersurface, since \(\Sigma\) has zero \((M-1)\)-dimensional measure.
Continuity therefore extends \(dh(V)\le0\) to that entire hypersurface.
A local projection \(\pi\) to it satisfies
\(\|x-\pi(x)\|\le C|h(x)|\).
The function \(a(x)=dh(x)V(x)\) is locally Lipschitz, so for \(h(x)\ge0\),
\(a(x)\le LC h(x)\).
Along a trajectory the positive part of \(h\) consequently satisfies the upper
right derivative inequality \(D^+h_+\le LC h_+\).
Gronwall's inequality forbids a crossing from \(h\le0\) to \(h>0\).
Every initial point of \(\operatorname{int}K\cap O\) outside the null hitting set
therefore remains in \(K\) until \(T\), because a first exit would occur in a
regular chart.  Such initial points are dense in \(K\).
Continuity of the flow and closedness of \(K\) finish the proof, since \(T\)
was arbitrary.
\end{proof}

\begin{proof}[Completion of the proof of Theorem~\ref{pos:finite}]
Necessity was proved in Lemmas~\ref{pos:finite-structure} and
\ref{pos:quadratic-witness}.  Conversely the rate conditions make \(A\) Metzler.
Its ambient flow \(e^{tA}\) is complete, and the strictly positive coefficient
orthant is forward invariant, since
\((e^{tA}f)_x\ge e^{tA_{x,x}}f_x>0\) for \(f\in O\).
Lemmas~\ref{pos:boundary} and \ref{pos:ordinary} verify the geometric hypotheses
of Lemma~\ref{pos:integration}.
At an ordinary boundary point, \(f(y)=0\) and \(\nabla f(y)=0\), so
\[
 dh_f(Af)=\sigma(Af)(y)=2\operatorname{tr}(Q(y)H)\le0.
\]
The normalization of \(H\) changes only a positive factor.
Lemma~\ref{pos:integration} now proves preservation of \(K\).
Invertibility of the finite matrix exponential excludes zero outputs from
nonzero inputs.  This proves sufficiency.
\end{proof}

\begin{remark}[A compact certificate]\label{pos:compact}
Write \(Q=Q^{[1]}+Q^{[2]}+Q^{[3]}\) by homogeneous degree and put
\(\widehat Q(r,y)=r^2Q^{[1]}(y)+rQ^{[2]}(y)+Q^{[3]}(y)\).
Condition~\eqref{pos:contact-test} is equivalent to
\(\operatorname{tr}(\widehat Q(r,y)H)\le0\) under
\eqref{pos:contact-domain}, \(r\ge0\), and \(r^2+\|y\|^2=1\).
For \(r>0\), this is the identity \(\widehat Q(r,y)=r^3Q(y/r)\);
the face \(r=0\) follows by continuity.  The set is compact, since
\(H\ge0\) entrywise and \(e^{\mathsf T}He=1\) bound every matrix entry.
The number of auxiliary variables is \(O(d^2)\), independently of the numerical
capacities.  This does not remove the capacities from the actual rates.
\end{remark}

\subsection{Mixed finite and infinite capacities}

Let \(F=\{i:\kappa_i<\infty\}\), \(U=\{i:\kappa_i=\infty\}\), and
\(J=\{i:\kappa_i=1\}\).  Let \(\mathcal{LP}_{+,\kappa}\) denote the positive
Laguerre--P\'olya class supported on \(\Omega_\kappa\).

\begin{theorem}[Mixed-capacity classification]\label{pos:mixed}
Let \(T_t\) be a coefficientwise-positive \(C_0\)-semigroup of bounded operators
on \(\ell^1(\Omega_\kappa)\) with generator \(A\).
Preservation of \(\mathcal{LP}_{+,\kappa}\) at every time is equivalent to
preservation for every allowed stable polynomial input, and to the following
structural assertion together with \eqref{pos:contact-test}:
all allowed polynomials belong to \(D(A)\), the only off-diagonal entries are
\begin{equation}\label{pos:mixed-rates}
\begin{aligned}
 A_{x-e_i,x}&=x_i\left[\delta_i+\alpha_i(x_i-1)+\sum_{j\ne i}a_{ij}x_j\right],\\
 A_{x+e_i,x}&=(\kappa_i-x_i)\Bigl[b_i+\chi_i(\kappa_i-x_i-1)\\
 &\hspace{3em}+\sum_{j\in F\setminus\{i\}}g_{ij}(\kappa_j-x_j)\Bigr] &&(i\in F),\\
 A_{x+e_j,x}&=\lambda_j &&(j\in U),\\
 A_{x-e_i+e_j,x}&=c_{ij}x_i(\kappa_j-x_j) &&(j\in F,\ i\ne j),\\
 A_{x-e_i+e_j,x}&=r_{ij}x_i &&(j\in U,\ i\ne j).
\end{aligned}
\end{equation}
Every displayed off-diagonal parameter is nonnegative.  The diagonal is
\eqref{pos:diagonal-rate}.  In \eqref{pos:principal-matrix} and
\eqref{pos:contact-domain} use the conventions
\[
 \begin{gathered}
 \alpha_i=\tau_i=0\ (i\in J),\qquad
 \chi_i=0\ (i\in U\cup J),\\
 g_{ij}=0\ \text{unless }i,j\in F,\qquad
 c_{ij}=0\ \text{unless }j\in F.
 \end{gathered}
\]
Entries with targets outside \(\Omega_\kappa\) are omitted.
The existence criterion for the formal operator is separate and is given in
Theorem~\ref{ana:weighted-generation}.
\end{theorem}
\begin{proof}
Assume polynomial-input preservation.  Jensen approximation, positive
coefficient closure, and the strictly positive diagonal entries of \(T_t\),
proved in Theorem~\ref{pos:symbol}, extend it to all
\(\mathcal{LP}_{+,\kappa}\).
Corollary~\ref{ana:polynomial-domain} supplies the polynomial domain and the
local first variation needed for
Lemma~\ref{pos:order}.  The order is at most two.
The growth-reduction theorem, Theorem~\ref{ana:growth-reduction}, applies with
Lemma~\ref{pos:poisson}: it identifies the bounded part that strictly increases
the total occupation of the unrestricted sites as
\[
 R=\sum_{j\in U}\lambda_jz_j+
   \sum_{i\in F,\,j\in U}r_{ij}z_j\partial_i,
 \qquad \lambda_j,r_{ij}\ge0,
\]
and supplies the locally uniform first variation at arbitrary real points.
In particular, \(L=A-R\) cannot increase that unrestricted occupation.
These analytic statements precede the signed contacts used below.

For completeness, we give the remaining algebraic reduction.
For a net shift \(\sigma\), the rate is a reduced polynomial of total degree
at most two.  Boundary interpolation, exactly as in
\eqref{pos:boundary-factors}, gives
\begin{equation}\label{pos:mixed-boundaries}
 p_\sigma(x)=\prod_i(x_i)_{(-\sigma_i)_+}
 \prod_{i\in F}(\kappa_i-x_i)_{(\sigma_i)_+}\,h_\sigma(x),
 \quad
 \deg h_\sigma\le2-\sum_i(-\sigma_i)_+-\sum_{i\in F}(\sigma_i)_+.
\end{equation}
For unrestricted coordinates, vanishing on every integer point is a polynomial
identity; for finite ones, use reduced interpolation.  Thus at most two source
or finite-target boundary factors can occur.
Double finite births have constant times two hole factors and are excluded by
\(A1\): its simultaneous two-birth coefficient is nonnegative but the first
variation of the second logarithmic coefficient along a positive diagonal is
nonpositive.  Hence that coefficient is zero.

For every allowed multi-index \(\alpha\) of total degree two, apply the product-root
test to \(\prod_i(z_i+s_i)^{\alpha_i}\).  It gives
\[
 P_\alpha(-s_{\supp\alpha},z_{\mathrm{spectators}})\le0
\]
for every real spectator vector.  We classify all contributions to this complete
coefficient polynomial, including catalytic one-particle shifts.
Quadratic monomials in unrestricted destinations outside the source
have nonnegative coefficients, by \eqref{pos:mixed-boundaries} and positivity.
Set these spectators equal and let them tend to positive infinity.  Their
quadratic coefficients must vanish.  A remaining free destination enters linearly.
For two distinct source sites its coefficient has the form
\[
 d+t_i z_i+t_jz_j-b_i z_i^2-b_jz_j^2;
\]
for a repeated source it has the form \(d+t_i z_i\).
These are all possibilities because \eqref{pos:mixed-boundaries} leaves at most
two total factors: a two-to-one replacement gives the constant \(d\),
a source-dependent transfer gives a term \(t_i z_i\), and a simultaneous finite
birth gives its negative hole term \(-b_i z_i^2\).
The destination can have either real sign, so its coefficient vanishes identically
on the negative source orthant.  Coefficient comparison removes every such
replacement, every source-self or third-site catalytic transfer, and every
simultaneous finite birth and unrestricted transfer.  A pure double death is the
remaining nonnegative constant in \(P_\alpha\); sending the negative source roots
to zero removes it as well.

A transfer to a finite site already has its two boundary factors and is
\(c_{ij}x_i(\kappa_j-x_j)\), \(c_{ij}\ge0\).
After the preceding eliminations a transfer to an unrestricted site can only be
\(x_i(r_{ij}+h_{ij}x_j)\), with \(r_{ij},h_{ij}\ge0\).
The pair contact at \((z_i,z_j)=(-s,-t)\) contains the term
\(h_{ij}t^2\) in its second-order coefficient, while the other remaining terms
are at most linear in \(t\) for fixed \(s\), apart from a possible reverse
finite-target contribution \(-c_{ji}s^2\).
Letting \(t\to\infty\) forces \(h_{ij}=0\).
Finally a finite birth with unrestricted catalyst has rate
\(h(\kappa_i-x_i)x_\ell\), \(h\ge0\), and contributes
\(-h z_i^2z_\ell\) to the mixed principal coefficient.
At \((-s,-t)\), divide the contact inequality by \(t\) and let \(t\to\infty\).
Its leading dependence on \(s\) is \(hs^2\); the remaining terms have degree
at most one.  Letting \(s\to\infty\) forces \(h=0\).
The surviving shifts are precisely \eqref{pos:mixed-rates}.
The small-root, large-root, and boundary-state arguments in
Lemma~\ref{pos:finite-structure} give all their asserted signs, also when the
other occupancies are unrestricted.  The diagonal is the reduced polynomial
of degree at most two given by normal ordering.
Lemma~\ref{pos:quadratic-witness} proves necessity of the contact inequality.

Conversely, impose artificial capacities \(N_j\ge2\) at the unrestricted sites.
Keep the original diagonal, and replace immigration and transfer into those
sites by
\(\lambda_j(1-x_j/N_j)\) and
\(r_{ij}x_i(1-x_j/N_j)\), respectively.
The principal trace changes by
\[
 \operatorname{tr}(Q_N(y)H)-\operatorname{tr}(Q(y)H)
 =-\sum_{j\in U,\,i\ne j}\frac{r_{ij}}{N_j}y_j^2H_{ij}\le0.
\]
The immigration correction is first order and the binary set is unchanged.
Every approximant therefore satisfies Theorem~\ref{pos:finite}.
Theorem~\ref{ana:killed-approximation} proves convergence of these finite flows,
started from their coordinate projections, to the given semigroup; its proof
uses the bounded increasing part and weighted estimates, not an unbounded
similarity on ordinary \(\ell^1\).
Theorem~\ref{ana:mixed-sufficiency} and Lemma~\ref{pos:lp-closure} pass preservation
to all \(\mathcal{LP}_{+,\kappa}\), and exclude zero outputs.
This proves all implications.
\end{proof}

\subsection{One-site cones, conservation, and exact obstructions}

\begin{corollary}[One finite site]\label{pos:one-site}
For capacity \(K\ge2\), the exact positive stability-preserving cone is
\[
 Az^n=d_nz^{n-1}+u_nz^{n+1}+v_nz^n,
\]
where
\[
 d_n=n[\delta+\alpha(n-1)],\quad
 u_n=(K-n)[b+\chi(K-n-1)],\quad
 v_n=c+\mu n+\tau n(n-1),
\]
\(\delta,b,\alpha,\chi\ge0\), \(c,\mu\in\R\), and
\(\tau\le2\sqrt{\alpha\chi}\).
For \(K=1\), every Metzler \(2\times2\) matrix qualifies.
For every \(K\ge2\), the cone modulo \(I,z\partial_z\) is nonpolyhedral.
\end{corollary}
\begin{proof}
In one variable the normalized contact has \(H=1\), \(y\le0\).
The condition \(\alpha y+\tau y^2+\chi y^3\le0\) is therefore
\(\alpha-\tau s+\chi s^2\ge0\) for all \(s>0\).
With \(\alpha,\chi\ge0\), this is equivalent to
\(\tau\le2\sqrt{\alpha\chi}\), including either zero parameter by limits.
For \(K=1\), every nonzero affine polynomial with nonnegative coefficients is
stable, and every positive invertible matrix exponential preserves this class.
For nonpolyhedrality, for each \(u>0\) the functional
\[
 \alpha/u+u\chi-\tau+\delta+b
\]
is nonnegative on the quotient cone.  Equality forces \(\delta=b=0\),
\(\alpha=u^2\chi\), and \(\tau=2u\chi\), by the arithmetic--geometric mean
inequality and the cone inequality.  Its zero set is exactly the ray generated
by \((\alpha,\chi,\tau)=(u^2,1,2u)\).
These distinct exposed rays, for every \(u>0\), exclude polyhedrality.
\end{proof}

\begin{corollary}[Conserving mixed systems]\label{pos:conserving}
A finite-exit, particle-number-conserving Markov generator on \(\Omega_\kappa\)
preserves all allowed positive stable polynomials if and only if its only rates are
\[
 q(x,x-e_i+e_j)=
 \begin{cases}
 c_{ij}x_i(\kappa_j-x_j),&i,j\in F,\quad c_{ij}=c_{ji}\ge0,\\
 r_{ij}x_i,&j\in U,\quad r_{ij}\ge0,
 \end{cases}
\]
and there is no transfer from an unrestricted to a finite site.
The preservation quantifier ranges over the entire allowed polynomial space,
including mixtures of different total particle numbers.
\end{corollary}
\begin{proof}
Every fixed-total state space is finite, so the chain is nonexplosive.
Conservation and nonnegative off-diagonal entries remove all births and deaths
from Theorem~\ref{pos:mixed}.
For two finite sites the pair part of the principal coefficient is
\[
 (c_{ij}+c_{ji})z_i z_j-c_{ij}z_j^2-c_{ji}z_i^2.
\]
The contact at \((-s,-t)\) requires
\((t-s)(c_{ij}t-c_{ji}s)\ge0\) for every \(s,t>0\).
Taking \(t\) on either side of \(s\) forces \(c_{ij}=c_{ji}\).
For \(i\in U,j\in F\), a rate \(c_{ij}x_i(\kappa_j-x_j)\) gives the pair
coefficient \(c_{ij}(z_i z_j-z_j^2)\).  Taking \(z_j=-s,z_i=-t\) and
\(t\to\infty\) forces \(c_{ij}=0\).
No second-order term restricts independent transfer into an unrestricted site.

For sufficiency, a finite pair contributes
\begin{equation}\label{pos:exclusion-generator}
 G_{ij}^\kappa=(z_j-z_i)(\kappa_j\partial_i-\kappa_i\partial_j)
                  -(z_j-z_i)^2\partial_i\partial_j.
\end{equation}
Normalized polarization replaces each site by \(\kappa_i\) binary slots.
Summing the swap generators over every cross-block pair of slots intertwines
with \(G_{ij}^\kappa\): on a state with occupancies \(x_i,x_j\), exactly
\(x_i(\kappa_j-x_j)\) swaps move a particle from \(i\) to \(j\).
A binary swap \(\tau\) has flow
\[
 e^{t(\tau-I)}=\frac{1+e^{-2t}}2I+\frac{1-e^{-2t}}2\tau,
\]
which preserves stability by partial symmetrization \cite{BBL09}.
Polarization, specialization, and finite-dimensional Lie products prove
preservation for finite exclusion.  Independent transfers into unrestricted
sites have the affine-substitution flow from Theorem~\ref{pos:symbol}.
On each fixed total-degree cutoff all these operators act in finite dimension,
so a final Lie product combines them.
\end{proof}

\begin{corollary}[Reversibility and conserving time-one embeddings]\label{pos:embedding}
On a fully finite box put \(w(x)=\prod_i\binom{\kappa_i}{x_i}\) and
\(D=\diag(\sqrt{w(x)})\).
Every generator in Corollary~\ref{pos:conserving} is reversible with respect to
\(w\).  If its site graph is connected, its unique stationary law at total
population \(n\) is
\(w(x)/\binom{\sum_i\kappa_i}{n}\).
A stochastic, number-conserving operator \(T\) is the time-one map of a positive
stability-preserving Markov semigroup if and only if
\(S=D^{-1}TD\) is symmetric positive definite and
\[
 D(\log S)D^{-1}=\sum_{\{i,j\}}c_{ij}G_{ij}^\kappa,
 \qquad c_{ij}\ge0.
\]
The embedding generator in this class is unique.
\end{corollary}
\begin{proof}
For \(y=x-e_i+e_j\), the binomial identities give
\[
 \frac{w(y)}{w(x)}=
 \frac{x_i}{\kappa_i-x_i+1}\,
 \frac{\kappa_j-x_j}{x_j+1}.
\]
This is exactly detailed balance for the symmetric rates.  Connectivity permits
moving particles along site paths until any prescribed occupancy of the same
total is reached, so the fixed-total chain is irreducible.  The normalization
is the coefficient of \(s^n\) in \(\prod_i(1+s)^{\kappa_i}\).

If \(T=e^A\) with \(A\) Markov and positive, choose \(a\) so that \(B=A+aI\)
is entrywise nonnegative.  If an entry of \(e^A\) between different total layers
is zero, every term in \(e^{-a}\sum_{m\ge0}B^m/m!\) at that entry is zero,
in particular the corresponding off-diagonal entry of \(A\).
Thus time-one conservation forces generator conservation.  The previous
corollary applies, and \(D^{-1}AD\) is symmetric.  Its exponential is symmetric
positive definite, and its unique symmetric logarithm is \(\log S\).
This proves necessity and uniqueness.  Conversely the displayed logarithm is
a nonnegative combination of the conserving generators
\eqref{pos:exclusion-generator}; its exponential is \(T\) and the previous
corollary proves preservation.
\end{proof}

For two binary sites, \(T_p=(1-p)I+p\tau\) preserves stability for all
\(0\le p\le1\), but its conserving Markov embedding exists exactly for
\(0\le p<1/2\).  Indeed the antisymmetric eigenvalue is \(1-2p\), and the
embedding generator is
\(-\tfrac12\log(1-2p)(\tau-I)\).
Furthermore the one-particle rates of any finite conserving system obey every
cycle identity: the product of forward rates along a labeled cycle equals
the product along its reverse, since \(a_{ij}=c_{ij}\kappa_j\).
Pointwise limits preserve these identities.  The independent directed
three-cycle on unrestricted sites violates them.  Thus unrestricted conserving
systems need not be dilute limits of finite conserving systems on the same
labeled sites.  This does not apply to the killed approximants used in
Theorem~\ref{pos:mixed}.

\begin{corollary}[Support restriction for universal conserving mixing]\label{pos:support-mixing}
On a full binary box, suppose a universally positive stability-preserving
conserving Markov flow leaves a fixed-size family \(\mathcal B\) invariant
and is irreducible on it.  Then, for a partition into site-graph components
\(C_1,\ldots,C_r\) and integers \(k_1,\ldots,k_r\),
\[
 \mathcal B=\{S:|S\cap C_\ell|=k_\ell\ \text{for every }\ell\}.
\]
\end{corollary}
\begin{proof}
Corollary~\ref{pos:conserving} makes the flow symmetric exclusion on a site graph.
Every edge transposition must preserve \(\mathcal B\), since a positive rate
cannot leave an invariant family.  Edge transpositions generate the full
symmetric group of each connected component.  Their orbits are exactly the
sets with prescribed component occupancies.  Irreducibility makes
\(\mathcal B\) one such orbit.
\end{proof}

\begin{proposition}[Exact compensation for a capped transfer]\label{pos:compensation}
Let \(z\) have finite capacity \(K\ge1\), and let \(w\) be unrestricted.
Add to common coalescence of rate \(\gamma\ge0\) the transition
\((m,n)\to(m+1,n-1)\) at rate \(cn(K-m)\), where \(c\ge0\).
The resulting Markov flow preserves positive stability if and only if
\(c\le2\gamma\).
\end{proposition}
\begin{proof}
The additional differential operator is
\(cK(z-w)\partial_w+cz(w-z)\partial_z\partial_w\).
For a contact matrix \(H=\left(\begin{smallmatrix}h&a\\a&k\end{smallmatrix}\right)\)
and \(y=(z,w)\), common coalescence has principal trace
\(\gamma(e^{\mathsf T}Hy-y^{\mathsf T}Hy)\le0\).
At \(c=2\gamma\), the total trace simplifies exactly to
\[
 \gamma e^{\mathsf T}Hy-\gamma(h+2a)z^2-\gamma k w^2\le0.
\]
The endpoints \(c=0\) and \(c=2\gamma\) therefore pass every contact, and
linearity gives the entire interval.  Theorem~\ref{pos:mixed} proves sufficiency.
For necessity take \(h=k=0,a=1/2\), \(z=-s,w=-t\), \(s,t>0\).
Twice the trace is
\(-\gamma(s+t)-(2\gamma-c)st-cs^2\).
If \(c>2\gamma\), choose \(s>\gamma/(c-2\gamma)\) and then sufficiently large
\(t\); the trace is positive.  This is an admissible contact, including when
\(K=1\), and proves necessity.
For example \(K=1,\gamma=1,c=3\), \(f=(z+2)(w+15)\), and \(t=\log2\) give
\[
 T_tf=30+\frac{139}{8}z+\frac38w+\frac14zw,
 \qquad
 \frac{139}{8}\frac38-30\frac14=-\frac{63}{64}<0.
\]
The last inequality is the bivariate multiaffine stability discriminant.
\end{proof}

\begin{proposition}[Capacity and unchanged coordinates]\label{pos:spectators}
The operator \(A=(z-z^2)\partial_z^2\) preserves positive stability on
\(V_{(2,1)}\), with the binary coordinate \(w\) unchanged.
It fails on \(V_{(2,1,1)}\) with two unchanged binary coordinates.
Moreover there is no positive stability-preserving generator \(\widetilde A\)
on three binary variables \((u,v,w)\) satisfying the exact compression
\(D\widetilde A P=A\), where
\[
 Pz=(u+v)/2,\quad Pz^2=uv,\quad Pw=w,
 \qquad Dg(z,w)=g(z,z,w).
\]
\end{proposition}
\begin{proof}
For one unchanged coordinate, the contact matrix is
\(H=\left(\begin{smallmatrix}h&a\\a&0\end{smallmatrix}\right)\), \(h,a\ge0\).
If \(h>0\), the condition \(Hy\le0\) forces \(y_z\le0\): use the second
row when \(a>0\), and the first row when \(a=0\).
Thus \(\operatorname{tr}(Q(y)H)=h(y_z-y_z^2)\le0\); if \(h=0\) it is zero.
Theorem~\ref{pos:finite} proves preservation.
With two unchanged binary variables, the stable input
\((z+w_1)(z+w_2)\) evolves into
\[
 \theta z^2+z(w_1+w_2)+w_1w_2+(1-\theta)z,
 \qquad\theta=e^{-2t}.
\]
Identify \(w_1=w_2=w\) and specialize \(z=r\in(0,1)\).
Its discriminant in \(w\) is
\(-4(1-\theta)r(1-r)<0\) for \(t>0\), proving failure.

Suppose the asserted lift exists.  Average it with the swap of \(u,v\).
The positive preserving generator cone is convex, by finite-dimensional Lie
products, and this averaging preserves the compression identity.
The averaged lift commutes with the swap.  Since \(PD\) is the projection onto
its symmetric subspace and \(D\) kills the antisymmetric subspace, we have
\(D\widetilde A=D\widetilde A PD=AD\).
Evaluation at all ones makes \(\widetilde A\) Markov, because \(A\) is Markov.
Positivity and this lumping identity forbid any change of \(w\), any jump from
zero \(u,v\)-occupancy, or any occupancy change from a singleton.
From the double occupancy \(uv\), symmetry forces deaths to \(u\) and \(v\)
at rate one each.  The only remaining transitions are swaps of the singleton
states, at rates \(r_0,r_1\ge0\) according to the value of \(w\).

For \(F_s=(u+w)(v+s)\), direct application gives at
\((u,v,w)=(2s,-s,-2s)\)
\[
 (\widetilde A F_s)(2s,-s,-2s)
 =s+(4-3r_0-6r_1)s^2.
\]
Put \(c=3r_0+6r_1\) and \(s=1/[2(1+c)]\).  This value is
\((c+6)/[4(c+1)^2]>0\), although both nonnegative affine factors vanish at
the chosen real point.  Lemma~\ref{pos:root-contact} gives a contradiction.
\end{proof}

For unrestricted \(z\)-coordinates, the Markov flow
\eqref{pos:markov-generator} with \(\lambda=0\) preserves all joint positive
stable inputs after adjoining one unrestricted unchanged coordinate if and only
if \(\gamma=0\).  Sufficiency is affine substitution.  For necessity, apply
Corollary~\ref{pos:markov} in the enlarged system: the common quadratic rate
must vanish because the unchanged coordinate has zero second-degree action.
It is enough to test \((w+a\cdot z)^m\), since positive scaling and translation
of \(w\) supply all enlarged multinomial tests.
This quantifier is different from the one-binary-spectator assertion above.


\begin{example}[The semigroup assumption in multinomial testing]\label{pos:isolated-tests}
Multinomial tests do not characterize arbitrary positive operators.
Define a stochastic linear map from two variables to one by
\[
 K(z_1z_2)=(1+w^2)/2,\qquad K(z^\alpha)=w\quad(\alpha\ne(1,1)).
\]
For \(F=(c+az_1+bz_2)^m\), put \(M=(c+a+b)^m\),
\(U=[z_1z_2]F\).  Then
\(KF=U(1+w^2)/2+(M-U)w\).
For \(m\ge2\),
\[
 M\ge\binom m2c^{m-2}(a+b)^2
 \ge 2m(m-1)c^{m-2}ab=2U.
\]
For \(m<2\), \(U=0\).  Thus the quadratic discriminant is
\(M(M-2U)\ge0\), and every such output is a nonzero positive stable polynomial.
Nevertheless the stable monomial \(z_1z_2\) is sent to a polynomial vanishing at
\(w=i\).
Even adjoining all monomial tests does not suffice if zero outputs are allowed.
The projection retaining only
\(z_1z_2,z_1z_4,z_2z_3,z_3z_4\) sends every multinomial, for \(m\ge2\), to
\[
 m(m-1)c^{m-2}(a_1z_1+a_3z_3)(a_2z_2+a_4z_4),
\]
and sends every monomial to a stable monomial or zero.
It sends \((z_1+z_2)(z_3+z_4)\) to \(z_1z_4+z_2z_3\), which vanishes at
\(z_1=z_4=1+i\), \(z_2=z_3=-1+i\).
These examples concern isolated maps and do not contradict
Corollary~\ref{pos:markov}.
\end{example}


\begin{corollary}[A finite-degree witness under an order bound]\label{pos:finite-witness}
Suppose a particle-number-conserving Markov generator on unrestricted sites is
known to have differential order at most \(R<\infty\).  If it is not an
independent-particle migration generator, then some multinomial of degree at
most \(\max(2,R)\) fails stability preservation at arbitrarily small positive
times.
\end{corollary}
\begin{proof}
If every first variation at a multinomial root contact of degree
\(3,\ldots,R\) passed Lemma~\ref{pos:root-contact}, coefficient comparison as
in Lemma~\ref{pos:order} would remove all derivatives of order greater than two.
Degree conservation already supplies the output-degree bound.
If all degree-two square contacts also passed, the positivity and mass arguments
in Corollary~\ref{pos:markov} would give common coalescence and independent
migration.  Conservation removes coalescence.  This contradicts the hypothesis.
Thus one of the stated root-contact conditions fails; its local root argument
produces nonreal zeros for arbitrarily small positive times.
\end{proof}

For example, the positive conserving binary-consensus generator
\(A=(z_1-z_2)^2\partial_1\partial_2\) sends \((z_1+z_2)^2\) at time \(t\) to
\((2-\theta)(z_1^2+z_2^2)+2\theta z_1z_2\), where \(\theta=e^{-2t}\).
Specializing \(z_2=1\) gives discriminant \(16(\theta-1)<0\) at every
positive time.


\section{A degree-two ancillary test for positive semigroups}
\label{anc:section}


Preservation of stable inputs with nonnegative coefficients need not
persist after adjoining unchanged variables.  The exact gap can be
detected with only two additional polynomial degrees.  The mixed-capacity
classification of Section~\ref{pos:section} yields an explicit criterion.

Let \(\kappa\in(\N_{>0}\cup\{\infty\})^d\), and let \(T(t)\) be an existing
coefficientwise-positive \(C_0\)-semigroup on
\(X_\kappa=\ell^1(\Omega_\kappa)\), with generator \(A\).
For a finite list of ancillary capacities \(\lambda\), the inert extension
\(T(t)\otimes\Id_\lambda\) acts coefficientwise in the original variables.
In the next theorem preservation means preservation of the corresponding
positive Laguerre--P\'olya class; equivalently, it suffices to test the
allowed stable polynomial inputs with nonnegative coefficients.
No assertion of semigroup generation is inferred from a formal operator.

\begin{theorem}[The sharp ancillary test]\label{gen:ancilla}
The following conditions are equivalent.
\begin{enumerate}
\item \(T(t)\otimes\Id_{(1,1)}\) preserves positive stability for every
      \(t\geq0\).
\item \(T(t)\otimes\Id_{(2)}\) preserves positive stability for every
      \(t\geq0\).
\item Every inert extension by finitely many finite- or infinite-capacity
      coordinates preserves positive stability for every \(t\geq0\).
\item All allowed polynomials belong to \(D(A)\), the generator has the
      mixed-capacity rate form of Theorem~\ref{pos:mixed}, and its symmetric
      second-order matrix in \eqref{pos:principal-matrix} satisfies
      \begin{equation}\label{gen:eq-entrywise-cone}
       Q_{ij}(y)\leq0\qquad(y\in\R^d,\ 1\leq i,j\leq d).
      \end{equation}
\item The same domain and rate form hold, and the rate parameters satisfy
      \begin{equation}\label{gen:eq-positive-complete-rates}
       \alpha_i=\chi_i=0,\quad \tau_i\leq0,\qquad
       a_{ij}=g_{ij}=0,\qquad
       \theta_{ij}^2\leq4c_{ij}c_{ji}.
      \end{equation}
      The capacity conventions of Theorem~\ref{pos:mixed} remain in force;
      in particular \(\tau_i=0\) for binary sites and \(c_{ij}=0\) when
      the destination \(j\) has infinite capacity.
\end{enumerate}
The two ancillary degrees are necessary uniformly over capacity vectors:
one inert binary coordinate cannot replace the two in condition (1).
For finite \(\kappa\), these conditions are also equivalent to complete
preservation of arbitrary complex stable polynomials by \(T(t)\).
\end{theorem}

\begin{proof}
First note that
\(X_{\kappa\sqcup\lambda}=\ell^1(\Omega_\lambda;X_\kappa)\).
The inert extension is therefore a positive \(C_0\)-semigroup of norm
\(\|T(t)\|\).  If polynomials belong to \(D(A)\), each finite polynomial
in the enlarged space belongs to its generator domain, and its generator
is \(A\otimes\Id\) there.  This verifies the analytic hypotheses whenever
Theorem~\ref{pos:mixed} is applied below.  On a preserving side that theorem
itself supplies the polynomial domain.  Clearly (3) implies (1) and (2).

Suppose (1) holds.  Inputs independent of the ancillary variables give
ordinary positive preservation.  Theorem~\ref{pos:mixed} supplies the rate
form, and its application on the enlarged space gives the contact test
\eqref{pos:contact-test} with \(\widetilde Q=Q\oplus0\).
Let \(u,v\) denote the ancillary coordinates.  For distinct original
indices \(i,j\), fix any \(y\in\R^d\), and set
\[
 a=e_i+e_u,\qquad b=e_j+e_v,\qquad
 H=\frac{ab^\top+ba^\top}{8},\qquad
 \widetilde y=(y,-y_i,-y_j).
\]
Here \(e_i\) are coordinate vectors and \(e\) below is the all-ones vector.
Direct calculation gives
\begin{equation}\label{gen:eq-ancillary-contact}
 \begin{gathered}
 H_{ab}\geq0,\qquad e^\top He=1,\qquad H\widetilde y=0,\\
 (He)(He)^\top-H=\frac{(a-b)(a-b)^\top}{16}\succeq0.
 \end{gathered}
\end{equation}
All diagonals required to vanish at binary sites do vanish.  Thus this is
an admissible contact in \eqref{pos:contact-domain}, and
\[
 0\geq\operatorname{tr}(\widetilde Q(\widetilde y)H)=Q_{ij}(y)/4.
\]
For \(i=j\) use \(a=e_i+e_u\), \(b=e_i+e_v\).
The same identities hold; the only nonzero original diagonal of \(H\)
is at site \(i\), so the construction is allowed when \(\kappa_i\geq2\).
It again gives \(Q_{ii}(y)\leq0\).  For binary sites that coefficient is
zero by canonical convention.  Hence (1) implies (4).

For (2), again obtain the original rate form first.  Fix \(y\) and the
indices \(i,j\), permitting \(i=j\) only at a nonbinary site.  Choose
\(R\geq\max(y_i,y_j)\), and put \(c=R-y_i\), \(d=R-y_j\).
With the single capacity-two ancillary coordinate \(w\), set
\[
 a=e_i+e_w,\quad b=e_j+e_w,\quad
 H=(ab^\top+ba^\top)/8,\quad \widetilde y=(y,-R).
\]
The normalization and semidefinite identity in
\eqref{gen:eq-ancillary-contact} still hold, while
\[
 H\widetilde y=-(da+cb)/8\leq0,\qquad
 \widetilde y^\top H\widetilde y=cd/4\geq0.
\]
The diagonal in coordinate \(w\) is now allowed.  The positive stable
quadratic is a positive scalar multiple of
\((z_i+w+c)(z_j+w+d)\).  The contact inequality once more isolates
\(Q_{ij}(y)/4\), proving (2)$\Rightarrow$(4).

A univariate polynomial \(\alpha s+\tau s^2+\chi s^3\) is nonpositive
on \(\R\) exactly when \(\alpha=\chi=0\) and \(\tau\leq0\).
For the bivariate entry in \eqref{pos:principal-matrix}, simultaneous
scaling of its two variables to both signs of infinity kills the cubic
homogeneous part; scaling towards zero kills the linear part.  The remaining
quadratic is nonpositive exactly when
\[
 \begin{pmatrix}c_{ji}&-\theta_{ij}/2\\
                 -\theta_{ij}/2&c_{ij}\end{pmatrix}\succeq0.
\]
The nonnegativity of the off-diagonal rate parameters makes this condition
precisely \eqref{gen:eq-positive-complete-rates}.  This proves (4)$\Leftrightarrow$(5).

Assume (4) and adjoin any finite list of ancillary capacities.  The
extended rate form is unchanged, with all new rates zero.  Every contact
matrix \(H\) in \eqref{pos:contact-domain} is entrywise nonnegative, so
\[
 \operatorname{tr}((Q(y)\oplus0)H)\leq0.
\]
Theorem~\ref{pos:mixed}, with the tensor-domain observation above, gives
preservation.  This proves (4)$\Rightarrow$(3).
In finite dimensions, (4) is exactly the second-order sign condition of
Theorem~\ref{gen:main-real}: the full mixed differential coefficient is
\(2Q_{ij}\), and the diagonal coefficient is \(Q_{ii}\).
That theorem proves the final complete-preservation assertion and its converse.

For sharpness take an active capacity-two coordinate \(z\) and
\(A=(z-z^2)\partial_z^2\).  With one inert binary coordinate its contact
matrix has the form \(H=\left(\begin{smallmatrix}h&a\\a&0\end{smallmatrix}\right)\),
where \(h,a\geq0\).  If \(h>0\), the condition \(Hy\leq0\) forces
\(y_z\leq0\), using the second row when \(a>0\), and the first when
\(a=0\).  Consequently
\(\operatorname{tr}(Q(y)H)=h(y_z-y_z^2)\leq0\); for \(h=0\) it is zero.
The finite positive criterion proves one-binary preservation.
With two inert binary variables the stable input \((z+u)(z+v)\) evolves to
\[
 e^{-2t}z^2+z(u+v)+uv+(1-e^{-2t})z.
\]
Identifying \(u=v=w\) and specializing \(z=r\in(0,1)\) gives a real
quadratic in \(w\) of discriminant
\(-4(1-e^{-2t})r(1-r)<0\) for \(t>0\).  Stability is preserved by both
operations, so the evolved polynomial is not stable.  This proves the
uniform sharpness, without asserting that every fixed capacity vector
requires two ancillary degrees.
\end{proof}

The complete positive cone has a particularly small pairwise representation.
Apart from scalar and linear Euler terms, its one-site and directed
first-order terms are deaths, finite births, infinite immigration,
transfers to infinite sites, and the damping operators
\(-\eta_iE_i(E_i-1)\), where \(E_i=z_i\partial_i\) and \(\eta_i\geq0\).
A pair with one finite and one infinite site additionally permits the
saturated transfer into the finite site,
\(c_{ij}z_j\partial_i(\kappa_j-E_j)\) for \(i\in U\), \(j\in F\),
with \(c_{ij}\geq0\); its pair diagonal parameter is necessarily zero.
For two finite sites the remaining terms are exactly
\begin{equation}\label{gen:eq-positive-pair}
 c_{ij}z_j\partial_i(\kappa_j-E_j)
 +c_{ji}z_i\partial_j(\kappa_i-E_i)+\theta_{ij}E_iE_j,
 \qquad \theta_{ij}^2\leq4c_{ij}c_{ji}.
\end{equation}
Indeed expansion recovers \eqref{pos:principal-matrix} and the prescribed
first-order rates; all remaining diagonal coefficients are scalar and
linear Euler terms.  Each pair cone is linearly isomorphic to
\(\mathbb S_+^2\).  Spectral decomposition therefore expresses every pair
as the sum of at most two rank-one terms, whose parameters have the form
\(c_{ij}=u^2\), \(c_{ji}=v^2\), \(\theta_{ij}=\pm2uv\).
This is a representation of the positive coefficient-preserving cone.
The signed apolar construction for general complex preservation is treated
in the companion manuscript~\cite{WeiGenerators2026}.


\section{Domains, generation, and finite approximation}
\label{ana:section}

The algebraic restrictions on a stability-preserving generator do not by
themselves specify a semigroup on an infinite-dimensional space.  This section
resolves the domain and approximation questions for summable positive
coefficients and gives separate criteria for the existence of a semigroup.

\subsection{An automatic domain theorem}

Let \(E\) be countable and let \(r:E\to\N\) have finite fibres
\(E_k=r^{-1}(k)\).  Write \(e_x\) for the standard basis of \(\ell^1(E)\).
A nonnegative sequence has interval support if its positive entries form an
interval of integers, possibly a singleton or an unbounded interval.

\begin{theorem}[Automatic finite-support domain]\label{ana:domain}
Let \(T(t)\) be a positive strongly continuous semigroup on \(\ell^1(E)\),
with generator \(A\).  Suppose that, for each \(x\), and all sufficiently small
\(t>0\), the sequence
\[
 b_k(t)=\sum_{y\in E_k}(T(t)e_x)_y
\]
is log-concave and has interval support.  Then \(e_x\in D(A)\), and
\[
 \supp(Ae_x)\subseteq E_{r(x)-1}\cup E_{r(x)}\cup E_{r(x)+1},
 \qquad E_{-1}=\varnothing .
\]
In particular, every finitely supported vector belongs to \(D(A)\), and
\(A\) maps finitely supported vectors to finitely supported vectors.
\end{theorem}

We give the transition-function argument needed in the proof, to distinguish
existence of coordinate derivatives from an a priori finite-rate assumption.
The classical result is due to Kolmogorov; see \cite{Chung67,Kendall55}.

\begin{lemma}\label{ana:transition-derivatives}
Let \(S(t)\) be a standard countable substochastic transition kernel.  Then
every off-diagonal limit \(S_{yx}(t)/t\), as \(t\downarrow0\), exists and is
finite.  The limit \((1-S_{xx}(t))/t\) exists in \([0,\infty]\).
\end{lemma}

\begin{proof}
Transpose to the row convention and add an absorbing cemetery state.
We obtain a stochastic kernel \(P(t)\), with
\(P_{ij}(t)\to\delta_{ij}\) at zero.  Its entries are continuous: the
semigroup identity gives
\[
 |P_{ij}(t+h)-P_{ij}(t)|\le 1-P_{ii}(h),
\]
and the corresponding left estimate follows by writing \(t=(t-h)+h\).
Fix \(i\ne j\) and \(0<\varepsilon<1/3\).  Choose \(\delta_0>0\) so that
\(P_{ii}(s),P_{jj}(s)\ge1-\varepsilon\) for \(0\le s\le\delta_0\).
Consider the discrete chain with transition matrix \(P(h)\), started at
\(i\).  Let \(\tau\) be its first visit to \(j\), let
\(f_l=\Pr(\tau=l)\), and put
\(a_k=\Pr(X_k=i,\tau>k)\).
For \(kh\le\delta_0\), first-visit decomposition and
\(P_{ji}(s)\le\varepsilon\) give
\[
 P_{ii}(kh)
 =a_k+\sum_{l=1}^k f_lP_{ji}((k-l)h)\le a_k+\varepsilon .
\]
Consequently \(a_k\ge1-2\varepsilon\), and
\(f_{k+1}\ge a_kP_{ij}(h)\).  If \(nh\le\delta_0\), then
\[
 P_{ij}(nh)
 =\sum_{l=1}^n f_lP_{jj}((n-l)h)
 \ge n(1-\varepsilon)(1-2\varepsilon)P_{ij}(h)
 \ge n(1-3\varepsilon)P_{ij}(h).
\]
Fix \(0<\delta\le\delta_0\) and take \(n=\lfloor\delta/h\rfloor\).
Continuity implies
\[
 \limsup_{h\downarrow0}\frac{P_{ij}(h)}h
 \le \frac{P_{ij}(\delta)}{\delta(1-3\varepsilon)}<\infty .
\]
First let \(\delta\downarrow0\) along a sequence attaining the lower limit,
and then let \(\varepsilon\downarrow0\).  The off-diagonal assertion follows.

Subdivision of time, retaining only the path which stays at \(i\), gives
\(P_{ii}(t)>0\) for every \(t\).  Thus
\(\phi(t)=-\log P_{ii}(t)\) is continuous, nonnegative, and subadditive.
Writing \(\delta=nt+s\), \(0\le s<t\), yields
\(\phi(\delta)\le n\phi(t)+\phi(s)\).  Hence
\(\liminf_{t\downarrow0}\phi(t)/t\ge\phi(\delta)/\delta\).
Taking the supremum over \(\delta>0\) proves existence of the limit, possibly
infinite.  Since \((1-e^{-u})/u\to1\) at zero, this is also the asserted
diagonal limit.
\end{proof}

\begin{proof}[Proof of Theorem~\ref{ana:domain}]
Write \(p_{yx}(t)=(T(t)e_x)_y\).  Choose \(M,\omega\) with
\(\|T(t)\|\le Me^{\omega t}\), and choose
\(\lambda>\max(\omega,0)\).  Define the strictly positive bounded weights
\[
 w_x=\int_0^\infty e^{-\lambda s}\|T(s)e_x\|_1\,ds .
\]
No uniform positive lower bound on \(w_x\) is required.  Positivity,
Tonelli's theorem, and the semigroup law give
\[
 \sum_yw_yp_{yx}(t)
 =e^{\lambda t}
 \left(w_x-\int_0^t e^{-\lambda s}\|T(s)e_x\|_1\,ds\right).
\]
Therefore
\[
 S_{yx}(t)=e^{-\lambda t}\frac{w_y}{w_x}p_{yx}(t)
\]
is a standard substochastic kernel, with deficit
\[
 1-\sum_yS_{yx}(t)
 =w_x^{-1}\int_0^t e^{-\lambda s}\|T(s)e_x\|_1\,ds
 =t/w_x+o(t).
\]
Lemma~\ref{ana:transition-derivatives} supplies finite off-diagonal
derivatives \(q_{yx}=\lim_{t\downarrow0}p_{yx}(t)/t\).

Fix \(x\), put \(N=r(x)\), and use the \(b_k\) of the theorem.
Strong continuity gives \(b_N(t)\to1\).  Finiteness of adjacent fibres and
the coordinate derivatives imply \(b_{N-1}(t),b_{N+1}(t)=O(t)\).
Log-concavity and interval support give the geometric estimates
\[
 b_{N+j}(t)\le b_N(t)
 \left(\frac{b_{N+1}(t)}{b_N(t)}\right)^j,\qquad
 b_{N-j}(t)\le b_N(t)
 \left(\frac{b_{N-1}(t)}{b_N(t)}\right)^j .
\]
These remain valid when an adjacent coefficient vanishes, since the
corresponding tail then vanishes.  It follows that
\[
 \sum_{|k-N|\ge2}b_k(t)=O(t^2),\qquad
 \sum_{y\ne x}p_{yx}(t)=O(t).
\]
The second estimate also uses finiteness of the three central fibres.
Since \(w\) is bounded and \(w_x>0\), the off-diagonal column sum of \(S(t)\)
is \(O(t)\).  Its displayed deficit is \(O(t)\); hence
\(1-S_{xx}(t)=O(t)\).  The diagonal limit in the preceding lemma is thus
finite, and so is
\[
 q_{xx}=\lim_{t\downarrow0}\frac{p_{xx}(t)-1}{t}.
\]
On the finite union of the three central fibres, the difference quotient
converges coordinatewise and therefore in norm.  Outside that union its
\(\ell^1\) norm is \(O(t^2)/t=O(t)\).  This is precisely strong
differentiability of \(T(t)e_x\) at zero, with the stated support.
\end{proof}

\begin{corollary}\label{ana:polynomial-domain}
On any mixed-capacity space with finitely many coordinates, a positive
\(C_0\)-semigroup preserving \(\mathcal{LP}_{+,\kappa}\) automatically has all
polynomials in its generator domain.  Its action on each monomial has finite
support and changes total degree by at most one.
\end{corollary}

\begin{proof}
Use the proper grading \(r(x)=|x|\).  The diagonal specialization of
\(T(t)z^x\) belongs to the univariate positive Laguerre--P\'olya class.
Its coefficient sequence is log-concave with interval support: this follows
from negative-real-root factorization for its polynomial approximants and
passage to coefficient limits.  Theorem~\ref{ana:domain} applies.
\end{proof}

\subsection{Consistent cutoffs and weighted spaces}

Let \(\mathcal F\) denote the finitely supported vectors.  Suppose \(L\) is
a Metzler operator on \(\mathcal F\), meaning that its off-diagonal matrix
entries are nonnegative, and preserves
\[
 W_N=\operatorname{span}\{e_x:r(x)\le N\}.
\]
Suppose \(R\) is a bounded nonnegative operator whose entries raise the
grading by exactly one.  Write \(\rho=\|R\|_1\).

\begin{theorem}[Cutoff generation criterion]\label{ana:weighted-generation}
The operator \(L+R\) on \(\mathcal F\) has a positive \(C_0\)-semigroup
realization on \(\ell^1(E)\), with this action on \(\mathcal F\), if and only
if there are \(\tau>0\) and \(B<\infty\) such that
\[
 \sup_N\sup_{0\le t\le\tau}\|\exp(tL|_{W_N})\|_1\le B.
\]
When it exists, the realization is unique, its generator is the closure of
\(L+R\), and \(\mathcal F\) is a graph core.

For \(s\ge1\), put
\[
 X_s=\left\{f:\|f\|_s=\sum_x|f_x|s^{r(x)}<\infty\right\}.
\]
If \(\|U(t)\|_1\le Me^{\omega t}\) for the semigroup generated by the closure
of \(L\), with \(M\ge1\) and \(\omega\ge0\), then compatible semigroups satisfy
\[
 \|U(t)\|_{X_s}\le Me^{\omega t},\qquad
 \|R\|_{X_s}=s\rho,\qquad
 \|T^{(s)}(t)\|_{X_s}\le
 M\exp\bigl((\omega+Ms\rho)t\bigr).
\]
\end{theorem}

\begin{proof}
The finite flows are consistent by invariance of the \(W_N\).  The asserted
bound, followed by subdivision of time, bounds these flows on every compact
time interval.  They extend from \(\mathcal F\) to bounded operators \(U(t)\)
on \(\ell^1(E)\).  Density proves both the semigroup law and strong
continuity.  Finite-dimensional positivity proves positivity of the
extension.  The bounded positive perturbation \(R\) can be added by the
Dyson series, giving a positive semigroup \(T(t)\).

Conversely, subtract \(R\) from any proposed generator by bounded
perturbation.  Each \(W_N\) is contained in its domain and is invariant under
its action.  Uniqueness of classical solutions identifies the resulting
semigroup on \(W_N\) with \(\exp(tL|_{W_N})\).  In particular, it is positive
on \(\mathcal F\), hence positive everywhere, and its local boundedness
gives the necessary cutoff estimate.  It also proves uniqueness.

Let \(G\) denote the generator of the constructed semigroup \(U(t)\).
For a sufficiently large resolvent parameter \(\lambda\), its resolvent
preserves each \(W_N\), by the Laplace-transform formula.
Given \(f\in D(G)\), approximate
\((\lambda-G)f\) in \(\ell^1\) by vectors in \(\mathcal F\) and apply that
resolvent.  The resulting vectors lie in \(\mathcal F\) and converge to
\(f\) in the graph norm of \(G\).  Since \(G=L\) on \(\mathcal F\), this
proves \(G=\overline L\).  A bounded perturbation has the same domain and an
equivalent graph norm, proving the core assertion.

All entries of \(U(t)\) vanish unless the output grading is no larger than
the input grading.  Its weighted column sums are therefore bounded by the
unweighted column sums after dividing by the input weight.  This proves
the bound for \(U(t)\) on \(X_s\).  Density proves strong continuity there.
The identity \(\|R\|_{X_s}=s\rho\) follows directly from its one-step grading.
The \(k\)-th Dyson term has norm at most
\[
 M^{k+1}e^{\omega t}(s\rho)^k t^k/k!,
\]
and summation gives the stated bound.  These constructions are compatible
under \(X_s\hookrightarrow X_1\), term by term.
\end{proof}

\begin{corollary}\label{ana:column-sum}
If the column sums of \(L+R\) are bounded above, the realization in
Theorem~\ref{ana:weighted-generation} exists.
\end{corollary}
\begin{proof}
The column sums of each finite restriction of \(L\) are no larger than those
of \(L+R\).  The positive matrix exponential has norm at most \(e^{ct}\)
when its column sums are at most \(c\).  Apply the cutoff criterion.
\end{proof}

\subsection{Growth reduction and the legitimate scaling limit}

For mixed capacities, write \(F\) and \(U\) for the finite and unrestricted
coordinates and use \(r(x)=|x_U|\).  All spaces \(W_N\) below are finite
dimensional.  The next proposition supplies the analytic bridge between
local normal-order tests and tests with arbitrary real spectators.

\begin{proposition}\label{ana:growth-reduction}
Let \(T(t)\) be a positive \(C_0\)-semigroup on \(X_\kappa=\ell^1(\Omega_\kappa)\)
preserving \(\mathcal{LP}_{+,\kappa}\).  Suppose its polynomial action has
already been reduced to canonical differential order at most two.  Then
every canonical coefficient has unrestricted-coordinate degree at most two.
The part that increases \(|x_U|\) is the bounded operator
\[
 R=\sum_{j\in U}\lambda_jz_j+
   \sum_{\substack{i\in F\\j\in U}}r_{ij}z_j\partial_i,
 \qquad \lambda_j,r_{ij}\ge0.
\]
For \(L=A-R\), the conclusions of
Theorem~\ref{ana:weighted-generation} hold.  In particular, for every
polynomial \(f\), the expansion
\(T(t)f=f+tAf+O(t^2)\) holds locally uniformly on every bounded polydisc.
\end{proposition}

\begin{proof}
Corollary~\ref{ana:polynomial-domain} shows that all canonical coefficients
are polynomials: the triangular normal-order recursion involves only
images of finitely many monomials.  Since there are finitely many
coefficients after order reduction, and finite-coordinate exponents are
bounded, there is a constant \(C\) such that
\[
 \|Az^{(m,x_U)}\|_1\le C(1+|x_U|)^2.
\]
Thus the Taylor truncations of \(z_F^m e^{a\cdot z_U}\), \(a\ge0\),
converge in graph norm.  Closedness of \(A\) puts this exponential in its
domain and identifies its derivative by the differential expression.

Specialize finite variables to \(b>0\); this specialization is bounded
because their degrees remain bounded.  Dividing the first variation by
the initial exponential gives
\[
 \sum_{|\alpha|\le2}(m)_{\alpha_F}b^{-\alpha_F}
       a^{\alpha_U}P_\alpha(b,z_U).
\]
Lemma~\ref{pos:poisson} makes this polynomial quadratic in
\(z_U\).  Comparing powers of \(a\), and then the linearly independent
falling factorials on the finite \(m\)-grid, shows
\(\deg_{z_U}P_\alpha\le2\).

Represent the operator as a matrix on the finite-coordinate state space:
\[
 A_{nm}=\sum_{\substack{|p|\le2\\|q|\le2}}
              C_{nm;pq}z_U^q\partial_U^p .
\]
Terms increasing \(|x_U|\) have \((|p|,|q|)=(0,2),(1,2)\), or \((0,1)\).
The coefficients with \((0,2)\) are nonnegative, as coefficients of the
image of the input with no unrestricted particles.  Poisson tangency at
\(a=0\) makes their quadratic part nonpositive on the positive orthant;
therefore they vanish.

For \((1,2)\), fix \(p=e_l\).  If \(q_l>0\), its coefficient is a slope in
the nonnegative affine rate for adding one unrestricted particle.
Letting \(x_l\to\infty\) proves that the slope is nonnegative.  If \(q_l=0\),
the shift removes a particle at \(l\) and adds two elsewhere, and its rate
is precisely that coefficient times \(x_l\); the coefficient is again
nonnegative.  After the \((0,2)\) elimination, apply Poisson tangency at
\(a=\varepsilon v\), \(v>0\), divide its quadratic inequality by
\(\varepsilon\), and let \(\varepsilon\downarrow0\).  Every remaining
linear-in-\(a\), quadratic-in-\(z_U\) coefficient is nonnegative and their
sum is nonpositive.  They all vanish.  The increasing part is consequently
\[
 R=\sum_{j\in U}z_jK_j,
\]
with \(K_j\) a nonnegative matrix on the finite state space.  This proves
boundedness before any scaling is performed.

Subtracting this bounded operator gives \(L=A-R\).  Its finite restrictions
are Metzler and preserve \(W_N\); the necessity part of
Theorem~\ref{ana:weighted-generation} supplies the positive semigroup
\(U(t)\) and its compatible weighted realizations.  Put
\[
 S_s:X_s\longrightarrow X_1,\qquad
 (S_sf)(z_F,z_U)=f(z_F,sz_U).
\]
This is an isometric isomorphism between the indicated different spaces.
For
\[
 \widetilde U_s(t)=S_sU^{(s)}(t/s)S_s^{-1}
\]
we have \(\|\widetilde U_s(t)\|\le Me^{\omega t/s}\).
On a fixed \(W_N\), every generator entry is
\[
 s^{-1}s^{|y_U|-|x_U|}L_{yx}=O(s^{-1}).
\]
Thus \(\widetilde U_s(t)\to I\) strongly, uniformly on compact time
intervals, first on \(\mathcal F\), then by density on \(X_1\).
The perturbation of this rescaled semigroup is exactly \(R\), since
\(S_sRS_s^{-1}=sR\).  Its \(k\)-th Dyson term is bounded by
\[
 M^{k+1}e^{\omega t/s}\rho^k t^k/k!,
\]
and converges strongly to \(t^kR^k/k!\).  Dominated summation proves
\[
 S_sT^{(s)}(t/s)S_s^{-1}\longrightarrow e^{tR}
\]
strongly and locally uniformly in time.  Positive scaling and coefficient
closure of \(\mathcal{LP}_+\), Lemma~\ref{pos:lp-closure}, show that
\(e^{tR}\) preserves that class.

Set all unrestricted variables except \(z_j\) to zero.  We find that
\(e^{tz_jK_j}\) preserves positive stability.  Since this one-step-raising
operator is bounded on every \(X_s\), its exponential has a first variation
on every bounded polydisc.  Finite-site double contacts
with the real variable \(z_j\) as a spectator force every second-order
normal coefficient of \(K_j\) to vanish: its product with \(z_j\) has to be
nonpositive for both signs of \(z_j\).  The first-order box structure in
Proposition~\ref{gen:box-decomposition} then expresses \(K_j\) by independent deaths, capacity
births, and an affine diagonal.  Corollary~\ref{ana:polynomial-domain}
applied to \(1\) excludes the births: a birth in \(K_j1\) would create a
degree-two term in \(R1\).  Positivity gives
\[
 K_j=\lambda_j I+\sum_{i\in F}(r_{ij}\partial_i+m_{ij}z_i\partial_i),
 \qquad r_{ij}\ge0.
\]
Its remaining diagonal slopes need not yet have a sign.

If \(m=m_{ij}\ne0\), use the input \(c+u\) in the corresponding finite
variable \(u\), with \(c>\max(0,r_{ij}/m)\).  Up to a nonvanishing
exponential factor, the output is
\[
 c+e^{mtz}u+\frac{r_{ij}}m(e^{mtz}-1).
\]
For \(z=iv\), choose \(v>0\) with \(\sin(mtv)>0\).  Its root
\[
 u=-r_{ij}/m-(c-r_{ij}/m)e^{-mtz}
\]
then has positive imaginary part, a contradiction.  Hence every
\(m_{ij}=0\), proving the displayed form of \(R\), including
\(\lambda_j\ge0\).

Finally, the generator on every \(X_s\) agrees with \(A\) on polynomials.
If \(f\in W_N\), then \(A^2f\in W_{N+2}\).  Twice integrating the
semigroup equation and using the weighted estimate gives, for
\(0\le t\le\tau\),
\[
 \|T(t)f-f-tAf\|_s
 \le \frac{t^2}{2}M s^{N+2}
        e^{(\omega+Ms\rho)\tau}\|A^2f\|_1 .
\]
Finite-coordinate radii add only a fixed factor.  This proves the final
assertion on arbitrary bounded polydiscs.
\end{proof}

\begin{remark}
The order reduction used in the preceding proposition needs only the
first variation on a small polydisc, which follows directly from the
\(\ell^1\) generator derivative.  Negative collision points can be chosen
inside that polydisc; polynomial identity then gives the global
normal-order conclusion.  Arbitrary-radius first variations are proved
only after the bounded growth decomposition.  Thus no weighted
regularity is presupposed in deriving that decomposition.
\end{remark}

\subsection{Finite approximation with unchanged diagonals}

\begin{theorem}\label{ana:killed-approximation}
Let a mixed-capacity operator have the rate form and quadratic-contact
inequality of the positive classification, and suppose it has the
realization in Theorem~\ref{ana:weighted-generation}.  Impose artificial
capacities \(N_j\ge2\) at unrestricted sites.  Keep its diagonal entries
unchanged and replace the rates for immigration and transfer into such a
site by
\[
 \lambda_j(1-x_j/N_j),\qquad
 r_{ij}x_i(1-x_j/N_j),
\]
respectively.  Denote the finite operator by \(A_N\), and the coordinate
projection onto its box by \(P_N\).  Then every \(A_N\) satisfies the finite
positive stability criterion.  As all \(N_j\to\infty\), the following
convergence is strong on \(X_\kappa\), uniformly on compact time intervals:
\[
 e^{tA_N}P_N\longrightarrow T(t).
\]
\end{theorem}

\begin{proof}
In the convention that the second-order expression is
\(\sum_{i,j}Q_{ij}\partial_i\partial_j\) with \(Q\) symmetric, the corrected
principal matrix is
\[
 Q_N(y)=Q(y)-
 \frac12\sum_{\substack{j\in U\\i\ne j}}
 \frac{r_{ij}}{N_j}y_j^2(E_{ij}+E_{ji}).
\]
For each admissible contact matrix \(H\), whose entries are nonnegative,
\[
 \Tr(Q_N(y)H)=\Tr(Q(y)H)
       -\sum_{\substack{j\in U\\i\ne j}}
                 \frac{r_{ij}}{N_j}y_j^2H_{ij}\le0.
\]
The immigration correction is first order, and the binary sites are
unchanged because \(N_j\ge2\).  The finite contact theorem applies.

Split \(A_N=L_N+R_N\) by unrestricted-degree increase.  The increasing
parts satisfy \(\|R_N\|\le\rho\) and converge strongly to \(R\).
For a nonnegative input in the finite box, variation of constants gives
\[
 U(t)f-e^{tL_N}f
 =\int_0^tU(t-s)(L-L_N)e^{sL_N}f\,ds\ge0.
\]
The finite trajectory is polynomial-valued and belongs to \(D(L)\).
The integrand is nonnegative because only off-diagonal entries were
reduced; their diagonals agree exactly.  Therefore
\(\|e^{tL_N}P_N\|\le Me^{\omega t}\).

On every fixed \(W_m\), the matrices \(L_N\) converge to \(L\) once
\(N_j\ge m\).  Finite-dimensional exponential convergence and density,
with the preceding uniform bound, imply
\(e^{tL_N}P_N\to U(t)\) strongly and uniformly on compact intervals.
The Dyson expansions adding \(R_N\) have uniformly summable majorants
\[
 M^{k+1}e^{\omega t}\rho^k t^k/k!.
\]
Each fixed term converges to the corresponding term for \(U,R\);
summation proves the claim.  The projection \(P_N\), including at time
zero, is essential in this full-space statement.
\end{proof}

\begin{corollary}\label{ana:mixed-sufficiency}
The rate and contact conditions in the positive mixed-capacity
classification, together with the cutoff generation criterion, imply
preservation of \(\mathcal{LP}_{+,\kappa}\).
\end{corollary}
\begin{proof}
For stable polynomial input, sufficiently large \(P_N\) act as the
identity.  The finite outputs are stable and nonnegative.  Their
\(\ell^1\) limit belongs to \(\mathcal{LP}_{+,\kappa}\) by
Lemma~\ref{pos:lp-closure}, normalizing at \(\mathbf1\) if the probability-measure
form of that theorem is used.  A nonzero positive input has nonzero
output: the diagonal entry \(T(t)_{xx}\) is positive by subdivision of
time and strong continuity.  Positive Jensen approximants converge in
\(\ell^1\) to every member of \(\mathcal{LP}_{+,\kappa}\).
Boundedness of \(T(t)\) and the same closure theorem complete the proof.
\end{proof}

\subsection{Further generation criteria and topology}

For comparison with other analytic topologies, let
\[
 Y_R=\left\{f=\sum_xf_xz^x:\sum_x|f_x|R^{|x|}<\infty\right\},
 \qquad \mathcal E_{\rm ent}=\bigcap_{R\ge1}Y_R .
\]
These coefficient seminorms give the usual compact-open topology on
entire functions in finitely many variables.  The map \(f(z)\mapsto f(Rz)\)
is an isometry \(Y_R\to Y_1\), rather than a bounded similarity of \(Y_1\)
when \(R>1\).  In the unrestricted positive canonical expression below,
this conjugation replaces
\[
 \lambda\mapsto R\lambda,\qquad b\mapsto b/R,\qquad
 \nu\mapsto\nu/R,
\]
and leaves \(c,M,\beta,\eta\) unchanged.  With coordinate weights \(R_i\),
the off-diagonal coefficient \(m_{ij}\) becomes \(m_{ij}R_j/R_i\).
These assertions follow by conjugating multiplication and differentiation
on each monomial.

\begin{proposition}\label{ana:entire-cutoff}
Let \(L\) be a degree-nonincreasing polynomial operator, and let its
consistent finite-degree flows be \(U_N(t)\).  They extend to a locally
equicontinuous \(C_0\)-semigroup on \(\mathcal E_{\rm ent}\) if and only if
for every \(R\ge1\) and \(\tau>0\) there are \(S\ge1\) and \(C<\infty\)
such that
\[
 \|U_N(t)p\|_{Y_R}\le C\|p\|_{Y_S}
 \quad(N\ge0,\ 0\le t\le\tau,\ p\in\mathcal P_N).
\]
Such an extension is unique.
\end{proposition}
\begin{proof}
Necessity is the definition of local equicontinuity, since finitely many
coefficient seminorms can be dominated by their largest radius.
For sufficiency, Taylor polynomials are dense in every seminorm.
The estimate makes their evolved sequences Cauchy, uniformly in
\(0\le t\le\tau\), in every target seminorm.  Their limits define
continuous operators satisfying the same estimate.  Density transfers
the semigroup law, and uniform approximation by a polynomial orbit
proves strong continuity.  The same density proves uniqueness.
\end{proof}

\begin{example}\label{ana:topology-warning}
The Euler flow \(e^{tE}f(z)=f(e^tz)\) acts continuously on all entire
functions, but for every fixed \(R\) and \(t>0\) its norm on \(Y_R\)
is infinite, as the degree-\(n\) monomial has norm ratio \(e^{tn}\).
Conversely, the positive stability-preserving polynomial flow
generated by \(z\partial_z^2\) cannot extend globally to all positive
entire inputs.  Indeed,
\[
 [z]e^{tz\partial_z^2}z^n=n!t^{n-1}\qquad(n\ge1).
\]
This follows by applying the lowering operator \(n-1\) times and dividing
by \((n-1)!\) in its exponential.  The coefficient of \(z\) in the evolved
Taylor truncations of \(e^{az}\), with \(a>0\), is therefore
\(\sum_{n=1}^N a^nt^{n-1}\), which diverges when \(at\ge1\).
Polynomial stability of this flow can be seen on any fixed degree bound
\(N\): it is the dilute limit of the one-site generators
\((z-z^2/K)\partial_z^2\), \(K\ge N\), which satisfy the finite positive
criterion.  Their restrictions to degree at most \(N\) converge as
\(K\to\infty\).  This example does not satisfy coefficient-\(\ell^1\)
generation and does not contradict Theorem~\ref{ana:unary-l1}.
\end{example}

For the unrestricted positive canonical operator, write
\[
 A=cI+\lambda\cdot z+
 \sum_i\left(b_i+\sum_jm_{ij}z_j\right)\partial_i
 +E D_\nu-\beta E(E-1)-\sum_i\eta_i z_i^2\partial_i^2,
\]
where \(E=\sum_i z_i\partial_i\), \(D_\nu=\sum_i\nu_i\partial_i\),
and the signs are those in the positive classification.  On the probability
simplex define
\[
 F(p)=\nu\cdot p-\beta-\sum_i\eta_i p_i^2.
\]

\begin{proposition}\label{ana:quadratic-growth}
Existence of a positive coefficient-\(\ell^1\) realization forces
\(F(p)\le0\) throughout the simplex.  The strict inequality
\(\max F<0\) is sufficient for such a realization.
\end{proposition}
\begin{proof}
For a finite set of basis states \(S\), let \(B\) be the principal
matrix of the generator on \(S\).  Positivity gives
\[
 P_ST(t)|_S\ge e^{tB}
\]
entrywise.  Indeed, the left side, denoted \(G(t)\), satisfies
\(G'(t)=G(t)B+R(t)\), where \(R(t)\ge0\), by differentiating on the
finite set of generator-domain columns.  Variation of constants proves
the comparison.

Choose \(x^{(N)}/N\to p\), with zero coordinates when \(p_i=0\), and retain
the finite states \(x^{(N)}-\alpha\), where \(0\le\alpha\le x^{(N)}\)
and \(|\alpha|\le K\).  For sufficiently large \(N\), this is the
fixed index set of such \(\alpha\) supported on \(\{i:p_i>0\}\).
Divide their principal matrix by \(N^2\).  It converges to the matrix with diagonal
\[
 -D(p),\qquad D(p)=\beta+\sum_i\eta_i p_i^2,
\]
and rates \(\nu_i p_i\) from \(\alpha\) to \(\alpha+e_i\) when
\(|\alpha|<K\).  All linear terms vanish in this scaling.  Its exponential
has column mass at \(\alpha=0\) equal to
\[
 e^{-sD(p)}\sum_{j=0}^K\frac{(s\nu\cdot p)^j}{j!}.
\]
If \(T(t)\) is locally bounded by \(M\), apply the comparison at
\(t=s/N^2\) and let \(N\to\infty\).  Then let \(K\to\infty\).
We obtain \(M\ge e^{sF(p)}\) for every \(s>0\), forcing \(F(p)\le0\).

Conversely, the column sum at \(x\) is
\[
 (Az^x)(\mathbf1)
 =|x|^2 F(x/|x|)+O(|x|)+O(1),
\]
with a uniform remainder.  If \(F\le-\varepsilon\), these column sums
are bounded above.  Corollary~\ref{ana:column-sum} applies.
\end{proof}

\begin{proposition}[Exact unary generation criterion]\label{ana:unary-l1}
For
\[
 A=c+\lambda z+(b+mz)\partial_z+
             \nu z\partial_z^2-\chi z^2\partial_z^2,
 \qquad \lambda,b,\nu,\chi\ge0,\quad c,m\in\R,
\]
the positive coefficient-\(\ell^1\) realization exists if and only if
\[
 \chi>\nu
 \quad\hbox{or}\quad
 \chi=\nu\ \hbox{and}\ m+b\le0 .
\]
\end{proposition}
\begin{proof}
The terms \(cI+\lambda z\) are bounded.  Removing them from an existing
generator preserves generation; positivity of the resulting semigroup
follows from its Metzler restrictions to finite degree cutoffs and
density, rather than from an assertion that subtracting a positive
operator always preserves positivity.  Conversely, add \(\lambda z\)
by positive bounded perturbation and \(cI\) by multiplying the semigroup
by \(e^{ct}\).  We may set \(c=\lambda=0\).
The column sum is
\[
 (\nu-\chi)n(n-1)+(m+b)n.
\]
This proves sufficiency, while Proposition~\ref{ana:quadratic-growth}
excludes \(\chi<\nu\).

Suppose \(\chi=\nu\) and \(a=m+b>0\).  Then \(A=Q+aE\), where \(Q\)
is pure death with rates \(d_k=bk+\nu k(k-1)\).
On the finite cutoff, the elementary finite-state Feynman--Kac formula
(obtained by conditioning on the first holding time) gives
\[
 \|e^{tA}z^n\|_1
 =\mathbb E_n\exp\left(a\int_0^t N_s\,ds\right).
\]
If \(\nu>0\), fix \(t>0\) and choose \(K\) with
\(\sum_{k>K}d_k^{-1}<t/4\).
The independent exponential holding times \(\tau_k\) give
\[
 H_n=\sum_{k=K+1}^n\tau_k,\qquad
 Z_n=\sum_{k=K+1}^n k\tau_k.
\]
Markov's inequality yields \(\Pr(H_n>t)\le1/4\), while
\[
 \mathbb EZ_n=\nu^{-1}\log n+O(1),\qquad
 \operatorname{Var}Z_n=\sum_{k=K+1}^n k^2/d_k^2=O(1).
\]
For large \(n\), Chebyshev's inequality gives
\(\Pr(Z_n\ge(\log n)/(2\nu))\ge3/4\).
On an event of probability at least \(1/2\), both estimates hold, and
the integrated population up to time \(t\) is at least \(Z_n\).
Thus \(\|e^{tA}z^n\|_1\ge\tfrac12 n^{a/(2\nu)}\), contradicting boundedness.
If \(\nu=0\), direct solution instead gives
\[
 \|e^{tA}z^n\|_1
 =\left(e^{mt}+b\int_0^t e^{ms}\,ds\right)^n .
\]
The parenthesis is \(1+at+O(t^2)>1\) for small positive \(t\), giving
the same contradiction.
\end{proof}



\section*{Disclosure of artificial intelligence use}
Most of the mathematical content, proofs, and text in this article
were generated with substantial assistance from
OpenAI GPT-6 Pro in ChatGPT chat mode and GPT-6 Astra in Codex, both
using the ultra reasoning setting. These were the only AI tools used.
The author organized and directed the research through repeated
interactions, including questions, choices of research directions,
requests for extensions and critical review, and selection and
organization of the material. AI assistance was also used for
literature searches, proof audits, exact computational checks, and
manuscript revision. These checks do not constitute complete
proof-assistant verification, which is not claimed.
The author takes full responsibility for the accuracy, originality,
attribution, and conclusions of the article.


\makeatletter\def\@biblabel#1{#1.}\makeatother
\providecommand{\bysame}{\leavevmode\hbox to3em{\hrulefill}\thinspace}
\providecommand{\MR}{\relax\ifhmode\unskip\space\fi MR }
% \MRhref is called by the amsart/book/proc definition of \MR.
\providecommand{\MRhref}[2]{%
  \href{http://www.ams.org/mathscinet-getitem?mr=#1}{#2}
}
\providecommand{\href}[2]{#2}
\begin{thebibliography}{1}

\bibitem{BB09}
Julius Borcea and Petter Br{\"a}nd{\'e}n, \emph{The {Lee--Yang} and
  {P{\'o}lya--Schur} programs. {I}. {Linear} operators preserving stability},
  Invent. Math. \textbf{177} (2009), no.~3, 541--569.

\bibitem{BB10II}
\bysame, \emph{The {Lee--Yang} and {P{\'o}lya--Schur} programs. {II}. {Theory}
  of stable polynomials and applications}, Comm. Pure Appl. Math. \textbf{62}
  (2009), no.~12, 1595--1631.

\bibitem{BB10}
\bysame, \emph{Multivariate {P{\'o}lya--Schur} classification problems in the
  {Weyl} algebra}, Proc. Lond. Math. Soc. (3) \textbf{101} (2010), no.~1,
  73--104.

\bibitem{BBL09}
Julius Borcea, Petter Br{\"a}nd{\'e}n, and Thomas~M. Liggett, \emph{Negative
  dependence and the geometry of polynomials}, J. Amer. Math. Soc. \textbf{22}
  (2009), no.~2, 521--567.

\bibitem{Chung67}
Kai~Lai Chung, \emph{Markov chains with stationary transition probabilities},
  second ed., Grundlehren der mathematischen Wissenschaften, vol. 104,
  Springer-Verlag, Berlin--Heidelberg, 1967.

\bibitem{Kendall55}
David~G. Kendall, \emph{Some analytical properties of continuous stationary
  {Markov} transition functions}, Trans. Amer. Math. Soc. \textbf{78} (1955),
  no.~2, 529--540.

\bibitem{LVR10}
Thomas~M. Liggett and Alexander Vandenberg-Rodes, \emph{Stability on
  {$\{0,1,2,\ldots\}^{S}$}: Birth-death chains and particle systems}, Notions
  of Positivity and the Geometry of Polynomials (Petter Br{\"a}nd{\'e}n, Mikael
  Passare, and Mihai Putinar, eds.), Trends in Mathematics,
  Birkh{\"a}user/Springer Basel, Basel, 2011, pp.~311--329.

\bibitem{WeiGenerators2026}
Dongsheng Wei, \emph{Generators of stability-preserving semigroups, spectral
  gaps, and classical ground-state computation}, 2026, Unpublished companion
  manuscript.

\end{thebibliography}

\end{document}

